\documentclass[conference,10pt]{IEEEtran}

\usepackage{fontspec}
\usepackage{amsmath,amssymb}
\usepackage{booktabs}
\usepackage{tabularx}
\usepackage{array}
\usepackage{graphicx}
\usepackage{xcolor}
\usepackage{enumitem}
\usepackage{balance}
\usepackage{listings}
\usepackage{tikz}
\usetikzlibrary{arrows.meta,positioning,fit,shapes.geometric,backgrounds}
\usepackage{hyperref}
\usepackage{xurl}

\setmainfont{Songti SC}[AutoFakeBold=2,AutoFakeSlant=0.2]
\setsansfont{PingFang SC}
\setmonofont{Menlo}[Scale=0.82]
\XeTeXlinebreaklocale "zh"
\XeTeXlinebreakskip=0pt plus 1pt
\emergencystretch=1em

\hypersetup{
  colorlinks=true,
  linkcolor=black,
  citecolor=black,
  urlcolor=blue,
  breaklinks=true,
  pdfauthor={Scry 技术团队},
  pdftitle={Scry 技术架构与 Agent 实现文档}
}

\definecolor{scryblue}{HTML}{2563EB}
\definecolor{scrygreen}{HTML}{0F766E}
\definecolor{scryamber}{HTML}{B45309}
\definecolor{scryred}{HTML}{B91C1C}
\definecolor{scrygray}{HTML}{F3F4F6}
\definecolor{scryline}{HTML}{64748B}

\newcolumntype{Y}{>{\raggedright\arraybackslash}X}
\newcolumntype{C}[1]{>{\centering\arraybackslash}p{#1}}
\newcommand{\code}[1]{\texttt{#1}}
\newcommand{\filepath}[1]{\path{#1}}
\newcommand{\implemented}{\textcolor{scrygreen}{\textbf{已实现}}}

\setlist[itemize]{leftmargin=1.25em,itemsep=1pt,topsep=2pt}
\setlist[enumerate]{leftmargin=1.45em,itemsep=1pt,topsep=2pt}
\lstset{
  basicstyle=\ttfamily\scriptsize,
  breaklines=true,
  frame=single,
  rulecolor=\color{scryline},
  columns=fullflexible,
  showstringspaces=false,
  keywordstyle=\color{scryblue}\bfseries,
  commentstyle=\color{scrygreen},
  stringstyle=\color{scryamber}
}

\title{Scry 技术架构\\与 Agent 实现文档}

\author{\IEEEauthorblockN{Scry 技术团队}
\IEEEauthorblockA{文档版本：1.0.0\quad 适用架构：Scry Production Architecture\\
适用范围：平台服务、智能运维 Agent、Loop 自动化与安全控制}}

\begin{document}

\maketitle

\begin{abstract}
Scry 是面向研发、运维和平台工程团队的一体化可观测与智能运维系统。平台以 OpenTelemetry Collector 接收指标、日志和链路数据，以 ClickHouse 承担遥测分析存储，以 Go Query Service 提供统一查询、告警、仪表盘、组织权限和平台配置能力，并通过 React 前端形成从服务总览、链路追踪、日志检索、基础设施监控到故障处置的完整工作台。系统的核心特色是把可观测数据、证据治理、运行记忆和受控执行组成同一条 Agent 闭环：TopoMem 使用当前事故 Trace 的异常调用路径跨越服务词汇差异，在根因服务区域内按错误签名和记忆分块重排；证据链记录每次工具结果的来源、置信度、关联度与人工处置；MCP 工具协议把自然语言意图编译为确定性查询；Loop 则把诊断、安全决策、审批、执行和验证固化为可恢复流程。Scry Agent Service 基于 Mastra、AI SDK、Hono、MCP、PostgreSQL 与 WASM SQLite 实现模型适配、持久对话、TopoMem 检索、证据验证写回、A2UI 解释、Prompt Injection 防护、最小权限 SSH 和记忆自维护。本文按正式生产架构说明系统分层、完整功能、Agent 与 TopoMem 实现、数据模型、接口契约、安全控制、部署和运维流程，为开发、交付、实施和安全评审提供统一技术依据。
\end{abstract}

\begin{IEEEkeywords}
可观测性，智能运维，Agent，TopoMem，运行拓扑记忆，证据链，MCP，Loop，人工审批，RBAC
\end{IEEEkeywords}

\section{引言}

\subsection{文档定位}

本文是 Scry 的正式技术架构文档，描述生产版本的组件职责、服务边界、运行协议、数据结构和运维方式。文档同时承担四类用途：开发团队据此理解模块依赖和扩展点；平台团队据此完成部署、容量规划与日常运维；实施团队据此配置数据接入、模型、Loop 和 SSH 策略；安全团队据此审查身份继承、秘密管理、人工审批与审计闭环。

本文覆盖 Scry Web Console、Query Service、Agent Service、OpenTelemetry Collector、ClickHouse、Alertmanager、ZooKeeper、Schema Migrator 以及受管 SSH 主机。外部模型服务被视为可替换依赖，通过 OpenAI Responses 或 Anthropic Messages 协议接入。所有用户可见功能、后台调度、事件触发、审批恢复和审计链路均按正式交付状态描述。

\subsection{适用读者}

\begin{itemize}
  \item \textbf{后端与前端开发者：}关注接口、模块、状态流、数据模型和错误语义。
  \item \textbf{Agent 与自动化开发者：}关注模型配置、记忆、工具、A2UI、Loop、调度器和审批机制。
  \item \textbf{SRE 与平台工程师：}关注端口、依赖、健康检查、容量、备份、高可用和故障恢复。
  \item \textbf{实施与解决方案团队：}关注功能边界、数据接入、权限角色、工作流配置和主机纳管。
  \item \textbf{安全与合规人员：}关注 JWT、RBAC、密钥加密、SSH 主机指纹、命令白名单和审计日志。
\end{itemize}

\subsection{系统边界与术语}

Scry 将系统拆分为遥测数据面、平台控制面、智能执行面和交互呈现面。遥测数据面负责接收和存储 metrics、logs 与 traces；平台控制面负责查询、规则、仪表盘、组织和访问控制；智能执行面负责模型推理、工具、Loop、调度与审批；交互呈现面负责 Web 页面、实时流、A2UI 和运维操作。

本文使用以下术语：\textbf{Thread} 表示一个持久化 Agent 任务对话；\textbf{Run} 表示一次 Agent 或 Loop 执行；\textbf{Step} 表示模型步骤、工具调用或 Loop 节点；\textbf{Tool} 表示具备严格输入输出 Schema 的受控能力；\textbf{Loop} 表示可版本化、可调度、可暂停恢复的自动化流程；\textbf{A2UI} 表示由工具返回、前端按协议渲染的结构化界面数据；\textbf{TopoMem} 表示使用事故时运行拓扑约束候选服务区域、再以错误签名和内容相似度重排的运维记忆检索机制；\textbf{证据} 表示由遥测、平台对象、文档、Loop 或受控主机工具产生并具有来源与生命周期的事实材料。

\section{核心技术特色与运行主线}

\subsection{差异化技术主线}

Scry 不是在观测平台旁附加一个通用聊天入口，而是把 Agent 的每一步限制在可验证、可解释、可恢复的系统路径内。完整观测功能提供实时事实基础，TopoMem 提供跨事故复用的运维经验，证据链决定哪些材料可以进入结论和长期记忆，MCP 负责把模型意图收敛为受约束工具参数，Loop 负责把一次性诊断提升为可重复执行的组织流程。表~\ref{tab:core-characteristics} 按重要性列出系统主线。

\begin{table*}[t]
\caption{Scry 核心技术特色及其系统作用}
\label{tab:core-characteristics}
\centering
\small
\begin{tabularx}{\textwidth}{p{2.5cm}p{3.0cm}Y Y}
\toprule
\textbf{技术主线} & \textbf{解决的问题} & \textbf{实现机制} & \textbf{可验证输出} \\
\midrule
TopoMem 运行拓扑记忆 & 症状服务与根因服务词汇不一致，平面检索把 Agent 引向错误服务 & 从当前证据提取 Trace 边，枚举异常路径和多分支候选，在服务局部执行签名、词项、稠密与质量重排；无可靠路径时选择性回退 & 检索策略、症状服务、匹配路径、签名词、分块得分、回退原因和检索运行 ID \\
证据治理闭环 & 模型结论缺少来源，历史经验与当前事实混用 & 工具结果自动入证据库；计算置信度和关联度；支持接受、驳回、挂起、备注、选择和基于更新证据重答 & 证据 UUID、来源、状态、评分因子、关系、操作历史及回答修订记录 \\
确定性 MCP 工具协议 & 模型猜测底层 Query API、字段包装或资源 ID，造成不可控试探 & 自描述 JSON Schema、扁平高层参数、服务端校验、JWT 权限继承、结果预算和 A2UI 映射 & toolCallId、规范化输入、结构化输出、错误分类、耗时和审计事件 \\
安全 Loop 执行 & 复杂运维任务无法稳定重复，自动操作缺少审批、恢复和验证 & 顺序、条件、并行、循环、MCP、安全检查、审批、受控操作和验证节点统一编译为持久 Run & Workflow Version、步骤状态、检查点、审批快照、运行事件、结果与重试记录 \\
统一可观测数据基础 & Agent 缺少可信实时事实或绕过平台权限直接访问存储 & OpenTelemetry、ClickHouse 与 Query Service 统一指标、日志、链路、告警、仪表盘和基础设施查询 & 有界时间窗、服务和属性过滤、查询状态、原始记录、聚合结果与 Trace 关联 \\
\bottomrule
\end{tabularx}
\end{table*}

\subsection{证据、记忆与执行闭环}

图~\ref{fig:evidence-memory-loop} 给出 Scry 的主要运行闭环。TopoMem 位于“获得事故上下文”与“形成首个诊断假设”之间，作用不是替代当前遥测，而是利用已验证历史经验选择首个服务目标和工具方向。任何记忆命中都以不可信历史数据进入 Prompt，必须由本轮工具证据重新验证；只有被接受的证据或具有可用证据支撑的 Loop 结果才能写回长期记忆。该约束把论文提出的运行拓扑检索思想\cite{topomem}落实为带权限、证据、维护和审计边界的产品路径。

\begin{figure*}[t]
\centering
\begin{tikzpicture}[
  node distance=7mm and 6mm,
  every node/.style={font=\small},
  box/.style={draw=scryline,rounded corners=2pt,align=center,text width=26mm,minimum width=28mm,minimum height=9mm,fill=white},
  arrow/.style={-{Latex[length=2mm]},draw=scryline,thick},
  feedback/.style={-{Latex[length=2mm]},draw=scrygreen,thick,dashed}
]
  \node[box,fill=blue!7] (incident) {告警/提问/Loop 事件\\确定症状与时间窗};
  \node[box,fill=purple!7,right=of incident] (evidence) {当前证据上下文\\Trace、Log、Metric、对象};
  \node[box,fill=green!8,right=of evidence] (topomem) {TopoMem\\路径导航与签名重排};
  \node[box,fill=orange!8,right=of topomem] (agent) {Agent + MCP\\选择工具并验证假设};
  \node[box,fill=blue!7,right=of agent] (answer) {结论与处置\\证据引用和 A2UI};
  \node[box,fill=green!8,below=11mm of topomem] (writeback) {验证写回\\分块、来源、关系、质量};
  \node[box,fill=purple!7,left=of writeback] (maintenance) {自维护\\衰减、强化、合并、冲突};
  \node[box,fill=orange!8,right=of writeback] (loop) {Loop 固化\\审批、执行、复核、审计};
  \draw[arrow] (incident) -- (evidence);
  \draw[arrow] (evidence) -- (topomem);
  \draw[arrow] (topomem) -- (agent);
  \draw[arrow] (agent) -- (answer);
  \draw[feedback] (answer) |- (writeback);
  \draw[arrow] (writeback) -- (maintenance);
  \draw[arrow] (writeback) -- (loop);
  \draw[feedback] (maintenance) -| (topomem.south);
  \draw[feedback] (loop) -| (agent.south);
\end{tikzpicture}
\caption{Scry 证据、TopoMem、Agent 与 Loop 闭环。实线表示本轮处理，虚线表示验证后的经验回流。}
\label{fig:evidence-memory-loop}
\end{figure*}

闭环中的主次关系是明确的：当前遥测证据优先于历史记忆，服务端策略优先于模型工具意图，人工证据处置优先于自动评分，验证节点优先于操作节点的自述结果。TopoMem 提高首个目标选择质量，但不会把记忆直接提升为事故事实；Loop 提高执行复用程度，但不会绕过角色、审批或命令白名单。

\section{总体架构}

\subsection{设计目标}

Scry 的架构设计遵循以下原则。

\begin{itemize}
  \item \textbf{统一观测：}指标、日志、链路和异常共用查询与可视化入口，减少多套工具之间的上下文切换。
  \item \textbf{自托管数据面：}遥测数据保存在 ClickHouse，平台配置保存在 SQLite，Agent 记忆和运行状态保存在 PostgreSQL，主机端口默认绑定回环地址。
  \item \textbf{权限继承：}Agent 不直连 ClickHouse，而是携带当前用户 JWT 调用 Query Service，使原有 RBAC 继续生效。
  \item \textbf{证据驱动诊断：}Agent 指令要求先查询、后结论，并区分事实、推断和行动项；复杂诊断先发布计划。
  \item \textbf{拓扑引导记忆：}历史经验先按事故时异常调用路径定位服务区域，再按错误签名和分块内容排序；缺少可靠 Trace 时显式回退，避免伪造拓扑确定性。
  \item \textbf{验证后学习：}记忆写回必须绑定可用证据或 Loop 运行结果；已驳回、挂起和注入风险内容不能进入在线长期记忆。
  \item \textbf{受控操作：}工具参数由 Zod 约束，查询和输出设置上限，SSH 执行受主机、命令、超时和审批策略共同控制。
  \item \textbf{状态可恢复：}对话、工作流定义、运行事件、步骤状态和审批点持久化，服务重启后可以继续读取和恢复。
  \item \textbf{自动化可治理：}Loop 同时支持手动、Cron 和事件触发，所有入口都进入统一的运行记录、幂等和审计链路。
  \item \textbf{协议可替换：}模型端点、遥测协议和呈现协议通过适配层隔离，避免业务逻辑绑定单一供应商。
\end{itemize}

\subsection{分层视图}

图~\ref{fig:architecture} 给出系统分层。OpenTelemetry 定义厂商中立的遥测数据模型与传输规范\cite{otel}，ClickHouse 提供面向分析负载的列式存储\cite{clickhouse}。Scry 在数据接入和分析存储之上构建统一 Query Service，并以独立 Agent Service 承载智能执行、自动化调度和高风险操作控制。

\begin{figure*}[t]
\centering
\begin{tikzpicture}[
  node distance=5mm and 7mm,
  box/.style={draw=scryline,rounded corners=2pt,minimum height=8mm,align=center,font=\small,fill=white},
  layer/.style={draw=scryline,rounded corners=3pt,inner sep=6pt,fill=scrygray},
  arrow/.style={-{Latex[length=2mm]},thick,draw=scryline},
  dashedarrow/.style={-{Latex[length=2mm]},thick,dashed,draw=scryamber}
]
  \node[box,fill=blue!7] (browser) {React/Webpack 前端\\观测页面、AI 工作台、Loop 编辑器};
  \node[box,fill=green!7,right=25mm of browser] (external) {模型服务\\OpenAI Responses / Anthropic Messages};

  \node[box,fill=blue!7,below=10mm of browser] (query) {Go Query Service\\查询、RBAC、仪表盘、告警、配置};
  \node[box,fill=green!7,right=15mm of query] (agent) {Scry Agent Service\\Agent + TopoMem + Evidence + Loop};
  \node[box,fill=orange!8,right=15mm of agent] (ssh) {受管 SSH 主机\\白名单命令与审批};

  \node[box,fill=purple!7,below=10mm of query] (collector) {OpenTelemetry Collector\\OTLP gRPC/HTTP 接入与处理};
  \node[box,fill=purple!7,right=15mm of collector] (alert) {Alertmanager\\告警通知与状态};
  \node[box,fill=purple!7,right=15mm of alert] (libsql) {PostgreSQL + WASM SQLite\\对话、TopoMem、证据、Loop 与审计};

  \node[box,fill=cyan!8,below=10mm of collector] (ch) {ClickHouse\\metrics / logs / traces};
  \node[box,fill=cyan!8,right=15mm of ch] (sqlite) {Query Service SQLite\\用户、组织、仪表盘等对象};
  \node[box,fill=cyan!8,right=15mm of sqlite] (zk) {ZooKeeper\\ClickHouse 协调};

  \draw[arrow] (browser) -- node[left,font=\scriptsize]{REST/WebSocket} (query);
  \draw[arrow] (browser) -- node[above right,font=\scriptsize]{UI Message Stream} (agent);
  \draw[arrow] (agent) -- node[above,font=\scriptsize]{Bearer JWT} (query);
  \draw[arrow] (agent) -- node[right,font=\scriptsize]{模型 API} (external);
  \draw[dashedarrow] (agent) -- node[above,font=\scriptsize]{策略化执行} (ssh);
  \draw[arrow] (collector) -- (ch);
  \draw[arrow] (query) -- (ch);
  \draw[arrow] (query) -- (sqlite);
  \draw[arrow] (query) -- (alert);
  \draw[arrow] (agent) -- (libsql);
  \draw[arrow] (ch.south) to[bend right=18] node[below,font=\scriptsize]{协调} (zk.south);
\end{tikzpicture}
\caption{Scry 总体架构与信任边界。虚线表示必须经过策略和人工审批的 SSH 路径。}
\label{fig:architecture}
\end{figure*}

\subsection{服务职责}

\begin{table*}[t]
\caption{运行服务及职责}
\label{tab:services}
\centering
\small
\begin{tabularx}{\textwidth}{p{2.3cm}p{2.4cm}Y Y}
\toprule
\textbf{服务} & \textbf{技术/端口} & \textbf{主要职责} & \textbf{持久化与依赖} \\
\midrule
Frontend & React 18、Webpack 5；3301 & 观测页面、AI 对话、A2UI 渲染、Loop 画布、设置与管理界面 & 无服务端状态；依赖 Query Service 和 Agent Service \\
Query Service & Go 1.22、Gorilla Mux；8080/8085 & V3/V4 查询、服务和基础设施 API、仪表盘、告警规则、用户组织与 RBAC、集成、日志管道 & ClickHouse；本地 SQLite；Alertmanager \\
Agent Service & Node.js 22、Hono、Mastra、AI SDK、MCP；4111 & 模型适配、对话流、TopoMem、证据链、MCP 工具、安全决策、A2UI、scry-ops 策略与 Loop 编译 & PostgreSQL、WASM SQLite、Query Service、外部模型端点 \\
OTel Collector & OpenTelemetry Collector；4317/4318 & 接收 OTLP 指标、日志和链路，执行解析、资源补充、批处理与导出 & ClickHouse；Query Service 管理接口 \\
ClickHouse & 24.1.2 & 高吞吐遥测写入、时间范围过滤、聚合和高基数查询 & \code{clickhouse\_data}；依赖 ZooKeeper \\
Alertmanager & 0.23.7；9093 & 告警状态和通知通道执行 & \code{alertmanager\_data}；回调 Query Service 内部 API \\
Schema Migrator & 0.111.16 & 同步与异步创建/升级 ClickHouse 模式 & 一次性任务，位于 ClickHouse 就绪之后 \\
ZooKeeper & 3.7.1 & ClickHouse 协调与集群元数据 & \code{zookeeper\_data} \\
\bottomrule
\end{tabularx}
\end{table*}

\subsection{数据面与控制面}

系统存在四条相互隔离但可关联的数据路径。遥测写入路径从被观测应用经 OTLP 进入 Collector，再写入 ClickHouse；交互查询路径从浏览器或 Agent 进入 Query Service，由查询构建器生成 ClickHouse 查询并规范化返回；平台控制路径保存用户、组织、告警、仪表盘、集成和偏好；智能执行路径保存模型设置、Thread、UI Message、证据、TopoMem、Loop、Run、SSH 主机、策略和审计。将 Agent 放在 Query Service 之外而不是放入数据库内部，可以避免模型直接生成任意 SQL，也使 RBAC、查询格式、缓存、超时和结果大小限制集中在统一边界上。

四条路径通过稳定标识关联。遥测记录使用 service、trace、span、timestamp 和属性集合建立上下文；平台对象使用 user、org、dashboard、rule 等标识；智能执行使用 threadId、runId、stepId、toolCallId 和 workflowId。一次故障诊断可以从告警对象出发，经服务和时间窗口查询到链路与日志，再把证据、模型结论、审批和 SSH 审计关联到同一个 Run。

\subsection{四平面职责}

\begin{itemize}
  \item \textbf{接入平面：}负责 OTLP 接收、数据校验、资源属性补充、批量处理、模式适配和写入重试。
  \item \textbf{查询平面：}负责查询解析、权限校验、查询构建、缓存、数据库执行、后处理和流式进度。
  \item \textbf{智能平面：}负责 Prompt 组装、模型调用、工具选择、Loop 编排、调度、审批和结果持久化。
  \item \textbf{呈现平面：}负责页面路由、任务管理、实时消息、图表、表格、A2UI、配置和运维交互。
\end{itemize}

每个平面只访问其必要依赖。Collector 不读取 Agent 数据；Agent 不直连 ClickHouse；浏览器不持有数据库凭据；SSH 主机不需要访问平台数据库。该分离降低了组件被攻破后的横向移动范围，并使容量扩展可以按写入、查询、模型和前端访问分别进行。

\subsection{请求身份上下文}

浏览器请求携带 Scry JWT。Query Service 在路由级执行 Open、View、Edit、Admin 或 Self 权限判断；Agent Service 先验证登录状态和角色，再把原始 JWT 透传给 Query Service 工具调用。智能执行上下文可表示为

\begin{equation}
\begin{aligned}
  C_{req}=(&userId,orgId,role,token,\\
           &threadId,runId,mode).
\end{aligned}
\end{equation}

工具实现不自行扩大角色权限。Thread、Run 和审批操作依据 userId/orgId 校验所有权；管理模型、SSH 主机、工具策略和 Loop 定义需要 ADMIN；运行已启用的 Loop 和读取本人运行记录由登录用户完成。由此，模型只负责提出工具调用意图，最终执行权由服务端上下文、Schema 和策略共同决定。

\section{遥测、查询与存储架构}

\subsection{遥测写入路径}

Collector 对外暴露 OTLP gRPC 4317 和 OTLP HTTP 4318。接收器解析 Resource、Scope 和信号数据，处理器补充平台所需属性并按批次聚合，导出器把 metrics、logs 和 traces 写入对应 ClickHouse 数据库。Collector 在 ClickHouse、模式迁移器和 Query Service 健康后启动，写入过程可抽象为

\begin{equation}
  E \xrightarrow{\mathrm{OTLP}} C \xrightarrow{\mathrm{process}} D_{m,l,t},
\end{equation}

其中 $E$ 为埋点应用或采集 Agent，$C$ 为 Collector，$D_{m,l,t}$ 分别表示 ClickHouse 内的 metrics、logs 和 traces 数据集。模式迁移器先创建与版本对齐的本地表、分布式表、物化视图和索引结构，因此 Collector 不承担在线建表职责。批处理参数在吞吐、写入延迟和内存之间取平衡；导出失败进入重试队列，并通过 Collector 自身健康与日志暴露状态。

指标以 metric name、temporality、fingerprint、timestamp 和 value 为核心；日志以纳秒时间戳、severity、body、resource attributes、log attributes、traceId 和 spanId 为核心；链路以 traceId、spanId、parentSpanId、service、operation、duration、status、events 和 attributes 为核心。统一的 Trace/Span 标识使日志、异常和链路详情能够互相跳转。

\subsection{统一查询路径}

Query Service 的 \code{/api/v3/query\_range} 接收统一查询结构，分别调用 metrics、logs 和 traces 构建器。一次请求经过六个阶段：参数与时间范围校验、信号类型分派、查询语句构建、缓存判定、ClickHouse 执行、后处理与响应格式化。后处理支持表格化、缺口填充、公式、Having、Limit 和 Reduce 等操作，使前端面板与 Agent 工具获得一致结果。

查询层同时提供属性键/值自动补全、聚合属性、过滤候选、查询格式化、查询进度推送和日志 Live Tail。前端 Query Builder、仪表盘、告警规则和 Agent 的 \code{querySignals} 工具复用同一查询契约，不维护互相分叉的查询语义。WebSocket 或流式端点负责长查询进度，用户中止操作会取消前端消费并释放后续处理资源。

设请求时间窗口为 $[t_s,t_e]$，序列数为 $n$，每条序列的采样点为 $p_i$，则结果规模近似为

\begin{equation}
  S \propto \sum_{i=1}^{n}p_i, \qquad p_i \approx \frac{t_e-t_s}{\Delta_i}.
\end{equation}

Agent 系统提示要求限制时间范围和结果规模；服务端进一步将单次工具返回文本限制为 \code{SCRY\_AGENT\_MAX\_TOOL\_RESULT\_BYTES}，默认 200,000 字节。这一组合同时控制模型上下文消耗和意外的大查询传播。面向图表的查询通过步长和最大点数控制采样密度，面向表格的查询通过 Limit 控制行数，面向日志和链路的明细查询通过排序游标实现稳定翻页。

\subsection{存储职责分离}

ClickHouse 只保存遥测分析数据。Query Service 的 SQLite 保存用户、组织、仪表盘、规则和其他业务对象。Agent 的 WASM SQLite 文件保存模型设置、UI 消息、MCP 插件、SSH 主机和策略、安全审计、工作流定义和运行状态；Mastra Memory 使用 PostgreSQL 的独立 Schema。各存储不应混为同一备份单元：恢复可观测数据需要 ClickHouse，恢复平台配置需要 Query Service 数据库，恢复 Agent 对话和密钥密文需要 Agent SQLite、PostgreSQL 与同一个主密钥。

\subsection{数据生命周期与一致性}

遥测数据按照平台 TTL 设置在 ClickHouse 中自动淘汰，配置对象则长期保存并通过显式删除或版本升级管理。Thread 删除会级联移除 UI Message 和关联 Loop Run；Workflow 删除会先清理运行记录；Workflow 定义变更会把旧定义写入版本历史。API Key 和 SSH Secret 只保存密文，读取设置时仅返回 \code{hasApiKey/hasSecret} 状态。

实时流和持久化采用双分支处理：一个分支立即返回浏览器，另一个分支消费完整 UI Message Stream 并写入消息表。工作流步骤在开始、暂停、完成和失败时更新运行记录，因此页面刷新后仍能恢复任务历史。调度器为每次触发生成唯一 trigger key，结合 workflowId 和计划时间或事件 ID 做幂等检查，避免多副本或重复事件创建重复 Run。

\section{完整功能说明}

\subsection{观测与分析功能}

表~\ref{tab:functions-observability} 根据前端路由和 Query Service API 汇总面向用户的观测功能。

\begin{table*}[t]
\caption{观测、分析与告警功能}
\label{tab:functions-observability}
\centering
\small
\begin{tabularx}{\textwidth}{p{2.1cm}p{3.0cm}Y p{1.5cm}}
\toprule
\textbf{功能域} & \textbf{入口/对象} & \textbf{完整能力} & \textbf{状态} \\
\midrule
服务监控 & \code{/services}、服务详情 & 服务列表；RPS、错误率、P50/P90/P99 延迟；端点和顶层操作；基于 span 的服务指标 & \implemented \\
链路分析 & Traces Explorer、Trace Detail & 条件过滤、属性补全、保存视图、链路详情、跨度瀑布/火焰关系、Trace ID 定位 & \implemented \\
日志管理 & Logs Explorer、Live Logs、Pipelines & 日志查询与聚合、实时尾随、字段管理、保存视图、解析管道预览与版本化配置 & \implemented \\
指标查询 & Query Builder、仪表盘组件 & 指标元数据、PromQL/查询构建、公式和后处理、时间序列/表格可视化 & \implemented \\
基础设施 & Hosts 及 API & 主机、进程、Pod、PVC、Node、Namespace、Cluster、Deployment、DaemonSet、StatefulSet、Job 的列表和属性筛选 & \implemented \\
服务拓扑 & \code{/service-map} & 根据链路依赖构建服务关系图，用于识别上下游与故障传播路径 & \implemented \\
异常 & Exceptions、Error Detail & 异常分组、计数、详情、前后错误导航及与链路上下文关联 & \implemented \\
消息队列 & Kafka 页面与详情 & Consumer Lag、生产/消费吞吐、分区延迟、网络延迟和接入状态 & \implemented \\
仪表盘 & 列表、详情、Widget & 创建、读取、编辑、删除仪表盘；变量查询；多种面板与自定义查询 & \implemented \\
告警 & Rules、Overview、History & 规则 CRUD/测试、状态时间线、统计、Top Contributors、维护停机计划 & \implemented \\
通知通道 & Settings/Channels & 通道列表、创建、编辑、删除和测试；由 Alertmanager 执行通知 & \implemented \\
用量与状态 & Usage、Status & 遥测用量、版本、健康状态、功能标志和运行配置查看 & \implemented \\
\bottomrule
\end{tabularx}
\end{table*}

\subsection{平台管理功能}

\begin{table*}[t]
\caption{平台、团队与智能运维功能}
\label{tab:functions-platform}
\centering
\small
\begin{tabularx}{\textwidth}{p{1.9cm}Y p{2.3cm}}
\toprule
\textbf{功能域} & \textbf{能力} & \textbf{状态} \\
\midrule
接入引导 & 首次注册、组织初始化、应用/日志/基础设施/AWS/Azure 监控接入步骤 & \implemented \\
身份与团队 & 登录、邀请与批量邀请、用户和组织管理、密码重置、个人设置、ADMIN/EDITOR/VIEWER 权限 & \implemented \\
访问与摄取 & Access Token、Ingestion Key、数据保留时间（TTL）、Apdex 阈值、字段配置 & \implemented \\
集成 & 内置集成安装与配置；ClickHouse、云平台和消息队列等数据源接入 & \implemented \\
调试模式 & 管理员启停调试数据、生成和清理 OTLP 验证数据 & \implemented \\
AI 助手 & 多任务对话、诊断/规划/操作模式、附件、流式推理、工具结果、来源、重试与停止 & \implemented \\
证据工作台 & 工具证据卡片、置信度、关联度、关系、接受/驳回/挂起/备注/选择及基于更新证据重答 & \implemented \\
TopoMem & 运行拓扑检索、错误签名、分块预览、来源关系、检索解释、证据写回、自学习与自维护 & \implemented \\
Loop & 可视化节点编排、条件/并行/循环、审批暂停恢复、输入、标签、版本和运行记录 & \implemented \\
Loop 调度 & Cron、时区、事件类型、过滤器、幂等触发、自动运行和失败重试 & \implemented \\
AI 设置 & 模型供应商、官方/兼容端点、API Key、生成参数、系统补充指令、记忆窗口 & \implemented \\
SSH 安全 & 主机指纹、密码/私钥密文、连接测试、只读白名单、预配置命令、逐次审批、审计 & \implemented \\
\bottomrule
\end{tabularx}
\end{table*}

\subsection{应用性能监控}

服务页以 service name 为主索引汇总应用性能。列表展示请求速率、错误率和延迟分位数，并允许按环境、版本、集群及自定义 Resource Attribute 筛选。服务详情把时间范围作为共享上下文，同步驱动吞吐、错误、P50/P90/P99 延迟、端点列表和依赖关系。用户从异常曲线选择时间区间后，可以直接进入对应端点、链路或日志，避免重新输入过滤条件。

顶层操作用于把框架级 Span 聚合成可理解的入口操作，例如 HTTP Route、RPC Method 或消息消费处理器。系统计算每个操作的调用量、错误率、平均/分位延迟和对服务整体延迟的贡献。Apdex 阈值由管理员设置，Query Service 按满意、可容忍和失败区间计算服务体验分数。服务详情同时保留原始属性和值的跳转能力，便于定位某个版本、区域或客户分组下的性能退化。

\subsection{分布式链路}

Traces Explorer 支持时间范围、服务、操作、持续时间、状态、Span Kind 和任意 Resource/Span Attribute 的组合筛选。查询结果提供 Trace 列表、持续时间分布和聚合统计；保存视图将查询、列配置、排序和展示模式持久化。Trace Detail 读取完整跨度树，展示父子关系、事件、错误状态、属性和时间轴，并根据相对起止时间绘制瀑布或火焰关系。

链路详情支持跨信号关联。Span 中的 service、traceId、spanId 和时间范围可以生成上下文链接，跳转到同一请求附近的日志；异常记录可以回到产生异常的 Span；服务拓扑使用聚合后的调用边构建上下游关系。对于高延迟请求，用户可以按自身耗时、子跨度耗时和外部依赖耗时区分瓶颈所在层级。

\subsection{日志检索与处理}

Logs Explorer 提供结构化 Query Builder、自由文本条件、字段选择、排序、直方图和表格明细。属性自动补全区分 Resource、Scope 和 Log Attribute，并根据字段类型提供字符串、数字和布尔运算。Live Logs 使用流式接口持续接收最新日志，支持暂停、恢复和保持当前过滤条件。保存视图记录查询语句、时间偏移、列宽和展示字段，供团队复用。

日志解析管道用于把非结构化 Body 转换为可查询字段。管理员可以创建管道、配置匹配条件和解析算子、预览样本结果、保存新版本并控制启用顺序。字段管理允许把高频字段提升为索引字段，系统在更新前校验名称、类型和冲突。日志记录保留 traceId/spanId 时，界面直接提供链路跳转；错误日志还可以进入异常聚合视图。

\subsection{指标与统一查询}

指标查询支持时间序列、标量、表格和公式。Query Builder 选择 metric、聚合、过滤、Group By、函数和步长，后端根据 temporality 区分 Gauge、Cumulative 和 Delta 语义，并对 Counter 计算 rate/increase。多查询结果通过公式节点进行四则运算或比例计算，后处理完成缺口填充、单位转换、排序、Having 和 Limit。

统一 V3/V4 查询契约让相同查询对象可以用于仪表盘面板、告警评估、服务分析和 Agent 工具。属性补全接口只返回当前信号和时间范围内存在的键值；查询格式化接口用于展示最终执行语义；查询进度接口提供长查询状态。结果携带单位、时间戳、序列标签和元数据，前端根据面板类型选择 uPlot、表格、数值、直方图或其他可视化组件。

\subsection{基础设施监控}

基础设施模块覆盖 Host、Process、Pod、PVC、Node、Namespace、Cluster、Deployment、DaemonSet、StatefulSet 和 Job。每类资源都提供属性键、属性值和列表 API，列表请求包含时间窗口、过滤、排序和分页。主机视图重点展示 CPU、内存、磁盘、网络、负载和进程；Kubernetes 视图聚合资源请求、限制、实际使用、重启、状态和工作负载关系。

资源详情把基础设施指标与服务信号关联。例如 Pod 可以关联 service.name 和 deployment/environment，Node 可以查看其上运行的 Pod，Deployment 可以比较副本与版本。基础设施告警使用同一规则系统，Agent 也可以先列出服务和告警，再通过统一查询读取特定资源指标，形成从平台信号到主机巡检的闭环。

\subsection{服务拓扑与异常管理}

服务拓扑根据链路中的调用方向聚合节点和边。节点包含服务健康、请求量、错误率和延迟，边包含调用次数、错误率和跨服务延迟。时间范围或过滤条件变化会重新计算拓扑，用户可以从边直接进入调用链路，从节点进入服务详情。该视图适合识别共享依赖、扇入扇出、循环调用和故障传播路径。

异常模块按异常类型、消息、堆栈和服务上下文生成 Group。列表显示首次出现、最近出现、次数、受影响服务和趋势；详情展示代表样本、堆栈、属性、关联 Trace 和相邻错误。通过 Group ID、Error ID 和前后导航接口，用户可以在聚合视图与原始事件之间切换，并把异常条件保存为查询或告警。

\subsection{消息队列监控}

Kafka 监控覆盖 Consumer Lag、Producer/Consumer Throughput、Partition Latency 和 Network Latency。系统按 cluster、topic、partition、consumer group、client 和 broker 组织指标，既提供全局概览，也提供生产者、消费者和分区明细。接入页面验证采集状态，详情页支持按时间范围比较积压增长、消费速率和端到端延迟。

消息队列信号可以和服务、链路、日志及告警组合使用。当某个 consumer group 的 Lag 持续增长时，用户可以定位对应消费者服务，查看消费 Span、错误日志和实例资源；Agent 也可以把积压、服务错误率和主机负载放在同一诊断计划中。

\subsection{仪表盘与告警}

仪表盘支持列表、创建、复制、编辑、删除、变量和独立 Widget 页面。面板保存查询、公式、时间偏移、单位、阈值、图例、布局和显示选项；变量通过查询接口动态生成值，并注入所有引用该变量的面板。布局使用响应式网格，编辑状态和查看状态分离，避免误操作改变生产看板。

告警规则复用统一查询对象，并配置评估窗口、频率、阈值、恢复条件、标签、注释和通知通道。规则管理器定时执行查询，维护 Pending、Firing 和 Resolved 状态，并把状态发送给 Alertmanager。历史页面展示状态时间线、统计、Top Contributors 和总体转换；停机计划在指定窗口抑制通知但保留评估记录。通道管理支持创建、编辑、删除和发送测试消息。

\subsection{集成、身份与平台设置}

集成模块管理内置接入模板和已安装实例，覆盖数据库、基础设施、云服务和消息队列等场景。模板定义所需遥测、接入步骤、默认仪表盘和验证条件。安装完成后，系统根据实际采集信号显示健康状态，并把相关仪表盘和告警入口组织到同一集成详情中。

身份系统支持首个管理员注册、登录、邀请、批量邀请、邀请撤销、用户编辑、角色修改、密码重置和个人偏好。角色分为 ADMIN、EDITOR 和 VIEWER：ADMIN 管理组织、用户、AI、SSH 和 Loop；EDITOR 创建和修改观测对象；VIEWER 执行查询和查看已授权资源。组织设置管理名称、偏好、数据保留、Apdex、Ingestion Key 和 Access Token。

接入引导按应用监控、日志管理、基础设施、AWS 和 Azure 等场景提供步骤化配置。连接状态页面通过查询实际数据确认接入，而不是只检查配置存在。调试模式通过现有 OTLP 链路生成二十类诊断案例，覆盖健康基线、PostgreSQL 连接池与表膨胀、Valkey/Redis 连接上限与持久化延迟、Kafka 消费积压、TLS 证书过期、DNS 解析失败、gRPC 截止时间、HTTP 429、响应结构不兼容、内存泄漏与 OOM、CPU 和线程池饱和、磁盘空间与 I/O、配置漂移、服务崩溃、网络监听暴露、僵尸进程以及复合故障。

场景目录把用户症状、服务拓扑、根因节点、错误比例、异常事件、根服务日志、上游传播日志、干扰日志、业务指标、主机指标和消息指标组织为同一不可变定义。生成器按轻量、标准和高负载档位分别产生 2、6 和 18 条请求链路，并在同一数据集中保留正常依赖对照；父 Span 覆盖完整下游时长，根因异常落在实际失败节点，上游 Span 只记录传播后的 4xx/5xx 或超时状态。结构化日志携带 Trace ID 和对应服务 Span ID，故障指标同时写入十分钟前的健康基线点和当前故障点，因此 Agent 可以从跨信号相关性和趋势而非单一错误文本完成诊断。

布尔属性 \code{scry.debug} 用于标记模拟遥测。字符串属性 \code{scry.dataset} 的值为 \code{diagnostic-evaluation}，\code{scry.debug.case\_id} 使用不透明编号 \code{case-xxx}。这些属性不携带场景名称、根因服务或根因标签。公开目录 API 返回编号、名称、类别、难度、症状、拓扑、信号类型和故障比例；管理员真值 API 返回根服务、标准根因、关键证据、处置、验证点和不安全动作。端到端准确性工具据此自动生成数据、调用 Agent、解析结构化诊断，并按根服务、根因、证据、处置、验证和安全六个维度评分。清理流程只删除带调试标识的对象和自动告警规则，不影响真实遥测。

\section{Query Service 与前端实现}

\subsection{Go Query Service}

Query Service 是浏览器、Agent 和数据存储之间的领域网关。其 \code{APIHandler} 组合 Reader、DAO、Query Builder、Querier、Rule Manager、Alertmanager Manager、Integration Controller、日志解析管道和基础设施 Repository。V3 查询构建器按信号类型委派：指标使用 metrics v3，日志和链路在当前启动参数下使用 v4。HTTP 路由通过访问级别包装器区分 Open、View、Edit、Admin 和 Self 权限。

除读查询之外，Query Service 还管理如下控制对象：仪表盘及变量、Explorer 保存视图、告警规则和停机计划、通知通道、字段与摄取设置、用户与邀请、组织和角色、偏好、集成和日志管道。它同时在 8085 提供供 Alertmanager 等内部组件使用的私有 API。该服务承担“统一语义和统一授权”，因此 Agent 工具层只需暴露少量高层操作。

Query Service 启动时完成配置解析、日志初始化、数据库迁移、Reader/DAO 创建、查询缓存初始化、规则管理器启动、集成控制器注册和 HTTP Server 装配。公共端口 8080 面向浏览器和 Agent，私有端口 8085 面向 Alertmanager 等内部服务。健康通道持续报告服务状态，进程接收 SIGTERM 后依次停止 HTTP Server、规则任务和后台资源，保证容器编排可以进行优雅下线。

\subsection{查询执行管线}

统一查询请求包含 signal、query builder expression、时间范围、步长、变量、公式和后处理选项。HTTP Handler 先完成 JSON 解码、角色校验和业务参数校验，再把请求交给 Query Builder。Query Builder 不直接执行 SQL，而是根据 metrics/logs/traces 类型选择专用构建器，生成带参数约束的 ClickHouse 查询计划。Querier 使用缓存键生成器计算稳定 Key，对历史时间段优先读取缓存，对数据仍在流动的尾部窗口绕过缓存。

Reader 执行查询并返回列式结果。后处理层将不同信号的原始结果规范化为统一 Series/Table 结构，依次应用公式、缺口填充、Having、Limit、Reduce 和表格转换。Formatter 根据单位把原始数值转换为时间、字节、数据率、百分比、吞吐或布尔显示。最终响应携带 status、data 和错误类型；超时、取消、坏数据、无权限、冲突和内部错误映射到明确 HTTP 状态。

查询进度通过 WebSocket/流式 API 发布。每个长查询记录 query ID、阶段和完成比例，前端在用户取消或页面离开时停止订阅。Live Tail 使用独立流式路径，持续读取满足条件的新日志。自动补全请求设置短期缓存，减少每次键入都扫描高基数字段。

\subsection{规则、仪表盘与配置对象}

规则管理器把规则定义编译为周期任务。每个任务在评估时间点构造查询窗口，执行统一查询，计算阈值和恢复条件，并将标签、注释和状态传递给 Alertmanager。规则定义和运行任务分离，编辑规则会原子更新任务；停用、删除和维护窗口都有明确状态转换。历史 API 从状态记录中生成时间线、统计和贡献者聚合。

仪表盘、保存视图、通道、用户、组织和偏好通过 DAO 访问 Query Service SQLite。数据库迁移在服务启动阶段执行，确保对象结构与版本一致。写接口使用 Edit 或 Admin 权限，读接口使用 View 权限，用户自身信息和密码使用 Self 权限。所有配置对象通过稳定 ID 访问，删除操作对不存在对象采用幂等语义。

\subsection{前端应用架构}

前端基于 React 18、Redux、React Query、Ant Design、Grafana Data、uPlot/Chart.js 和 Webpack 5。页面按路由懒加载，主导航将 AI 助手和 Loop 放在服务、链路、日志、基础设施、仪表盘、告警等运维入口之前。Agent 页面使用 AI SDK 的 \code{useChat} 和 \code{DefaultChatTransport} 消费 UI Message Stream；Loop 编辑器使用 DnD Kit 实现节点拖放和排序。

Agent 工作区采用三栏结构：左侧管理任务线程，中间显示对话和输入，右侧检查执行、证据、产物和审批。消息渲染器处理文本、reasoning、来源、文件、工具步骤、A2UI、工作流事件、审批、暂停和 tripwire 等 AI SDK/Mastra 消息部件。表格型 A2UI 最多展示 20 行和 6 列，其他结构化结果回退为 JSON，避免未知结构导致界面崩溃。

应用路由统一声明页面组件、路径、私有状态和导航 Key。页面组件使用动态 Import 拆分 Bundle，首次进入某个功能域时再加载对应代码。身份信息、组织上下文和访问 JWT 保存在全局状态；服务查询、仪表盘和配置对象使用 React Query 缓存；编辑器局部状态使用组件 Hook 管理。Axios/Fetch 客户端统一注入 API Endpoint、Authorization 和错误处理。

观测页面复用 Query Builder、时间选择器、变量、图表、表格和保存视图组件。用户改变时间范围或过滤条件后，页面生成规范化查询对象，由 React Query 管理加载、缓存、失败和重试状态。图表 Hover、缩放和选区操作会更新共享时间范围，从服务曲线到链路、日志和异常的跳转携带相同上下文。

\subsection{Agent 工作区前端}

Agent Workspace 首次加载时读取 Thread、模型设置和已启用 Loop。如果用户没有 Thread，前端自动创建一个。Thread 按“今天”和“更早”分组，可搜索、切换和删除；标题从首条用户消息截取生成。选择 Thread 后，前端读取持久化 UI Message 并恢复工具、审批、来源和工作流事件，而不是只恢复纯文本。

输入区支持多行自动高度、Enter 发送、Shift+Enter 换行、文件选择、文件去重、单文件删除、停止执行和重新生成。选择 Loop 时，请求自动切换到 Workflow Stream Endpoint，并把输入字段序列化为执行消息；普通对话则发送 mode、Thread 和最后一条 UI Message。传输层依据审批恢复体中的 runId/resumeData 选择继续原 Agent Run 或启动新请求。

右侧 Inspector 把所有 Message Part 分类为执行、证据、产物和审批。工具 Part 展示状态、输入、输出与 A2UI；Source Part 展示 URL 或文档来源；Workflow Part 展示步骤完成比例、输入、输出和暂停载荷；Approval Part 提供拒绝与批准按钮。所有未知 data Part 仍可通过 JSON Details 展示，保证协议新增字段不会丢失。

\subsection{Loop 编辑器前端}

Loop Manager 由管理台和编辑器组成。管理台支持搜索、状态/触发方式筛选、创建、复制、删除和进入编辑。编辑器左侧是节点面板，中间是可排序画布，右侧是 Loop 或节点属性。用户可以拖入 Agent、服务、告警、仪表盘、审批、条件、并行和循环节点，并在嵌套编辑器中配置子步骤。

Loop 级配置包括名称、描述、启用状态、标签、Cron、时区、事件触发器和输入 Schema。节点级配置包括名称、提示词、审批文案、重试次数、条件路径、操作符、分支、并行分支、循环模式、最大次数、数据源路径和 foreach 并发。保存前前端保证至少一个节点；服务端继续校验 ID、数量、嵌套深度和类型完整性。运行历史显示状态、时间和输入，用户可以把历史输入重新带入 AI Workspace 执行。

\section{Agent 子系统实现}

\subsection{组件构成}

Agent Service 按职责由 Hono HTTP 边界、身份适配器、模型适配器、Mastra Agent 与会话 Memory、TopoMem 运行记忆引擎、证据引擎、MCP Server/Client、Scry/SSH 工具、安全策略、A2UI、Loop 编译器和双存储 Repository 组成。Mastra 提供 Agent 和 Workflow 抽象\cite{mastraagents,mastraworkflows}，AI SDK 提供模型供应商适配、工具调用和 UI Message Stream\cite{aisdktools}，Hono 提供轻量 HTTP 运行时\cite{hono}。PostgreSQL 保存 Mastra 会话 Memory 与运行状态，WASM SQLite 保存证据、TopoMem、Loop、插件、SSH 配置和审计，因此 Agent 在 LoongArch64 上不加载本地原生 SQLite 扩展。

HTTP 边界只负责认证、Schema、所有权和响应协议；Agent 负责推理和工具选择；Workflow 负责确定性控制流；Coordinator 负责把时间和事件转换为 Run；Repository 负责定义、状态、恢复点和审计。该划分防止模型逻辑直接操作数据库或调度器，也让手动、定时和事件入口共享同一个 Workflow Runtime。

\begin{figure}[t]
\centering
\begin{tikzpicture}[
  node distance=4mm,
  every node/.style={font=\scriptsize},
  box/.style={draw=scryline,rounded corners=2pt,align=center,minimum width=0.82\columnwidth,minimum height=6mm,fill=white},
  arrow/.style={-{Latex[length=1.7mm]},draw=scryline,thick}
]
  \node[box,fill=blue!7] (api) {Hono API：认证、Zod 校验、线程/设置/工作流/SSH 路由};
  \node[box,fill=green!7,below=of api] (runtime) {Runtime Factory：按用户、模型配置和模式创建 Mastra};
  \node[box,fill=purple!7,below=of runtime] (memory) {Evidence + TopoMem：运行拓扑、记忆分块、检索包与验证写回};
  \node[box,fill=green!7,below=of memory] (agent) {Mastra Agent：Instructions + Memory Packet + Model + Tools};
  \node[box,fill=orange!8,below=of agent] (tools) {MCP Tools / Security Policy / scry-ops / A2UI / Approval};
  \node[box,fill=purple!7,below=of tools] (store) {PostgreSQL + WASM SQLite：对话、记忆、证据、Loop、SSH 与审计};
  \draw[arrow] (api) -- node[right]{创建} (runtime);
  \draw[arrow] (runtime) -- node[right]{检索上下文} (memory);
  \draw[arrow] (memory) -- node[right]{装配} (agent);
  \draw[arrow] (agent) -- node[right]{调用} (tools);
  \draw[arrow] (tools) -- node[right]{读写状态} (store);
\end{tikzpicture}
\caption{Agent Service 内部组件。运行时按请求创建，持久化存储在请求之间共享。}
\label{fig:agent-components}
\end{figure}

\subsection{TopoMem 运行拓扑记忆引擎}

\subsubsection{技术定位与多跳词汇鸿沟}

传统会话 Memory 主要保存最近消息，平面 RAG 主要依据查询与文档的词项或语义相似度排序。微服务事故中，用户看到的症状和根因文档往往由不同团队描述：入口服务使用 HTTP 500、页面延迟、请求超时等词汇，而根因服务的运行手册使用连接池耗尽、锁等待、证书过期或消费位点等领域词汇。两端可能几乎没有共同词项，真正连接它们的边只存在于本次事故的 Trace。TopoMem 针对这一多跳词汇鸿沟，把事故时运行拓扑作为检索控制信号，而不是把调用链仅拼接成普通查询文本\cite{topomem}。

Scry 中 TopoMem 位于 Agent 创建阶段。Runtime Factory 读取当前 Thread 中状态可用的证据，过滤已驳回和已挂起项目，从证据正文提取 Trace 边、服务名和错误签名，再执行记忆检索。返回的 Memory Packet 在模型首次选择工具之前注入系统上下文，但被明确标记为“不可信历史经验”：它只能用于形成待验证假设和选择首个服务目标，不能替代本轮遥测证据，也不能被当作证据 ID 引用。

\begin{figure}[t]
\centering
\begin{tikzpicture}[
  node distance=3.5mm,
  every node/.style={font=\scriptsize},
  box/.style={draw=scryline,rounded corners=2pt,align=center,minimum width=0.83\columnwidth,minimum height=6mm,fill=white},
  arrow/.style={-{Latex[length=1.7mm]},draw=scryline,thick}
]
  \node[box,fill=blue!7] (input) {问题 + 可用证据：告警、Trace、Log、Metric、平台对象};
  \node[box,fill=purple!7,below=of input] (extract) {规则提取：症状服务、异常 CALLS 边、错误签名};
  \node[box,fill=green!7,below=of extract] (walk) {TopoMem-Walk / TopoMem-MP：最多 3 条异常路径};
  \node[box,fill=orange!8,below=of walk] (rank) {服务局部：分块词项 + 局部稠密 + 签名 + 质量重排};
  \node[box,fill=blue!7,below=of rank] (packet) {可解释 Memory Packet：路径、因子、分块、来源与回退原因};
  \draw[arrow] (input) -- (extract);
  \draw[arrow] (extract) -- (walk);
  \draw[arrow] (walk) -- (rank);
  \draw[arrow] (rank) -- (packet);
\end{tikzpicture}
\caption{TopoMem 在线检索管线。拓扑决定优先服务区域，文本信号决定区域内记录顺序。}
\label{fig:topomem-pipeline}
\end{figure}

\subsubsection{运行时拓扑提取与路径导航}

TopoMem 从两类结构提取有向边。第一类是具有 \code{spanId}、\code{parentSpanId} 和 \code{service.name} 的 Span 集合，父子 Span 跨服务时形成调用边；第二类是工具结果中显式的 \code{fromService/toService}、\code{upstream/downstream} 或 \code{caller/callee} 结构。边记录调用次数 $c_e$ 和错误率 $r_e$，只有 $r_e\geq0.1$ 的边进入异常子图。症状服务优先取用户问题中出现的服务名，其次取告警或服务证据中的首个服务，最后回退到第一条运行边的起点。

从症状服务开始，路径枚举对每个节点最多展开 3 条出边，单条路径最多包含 6 个服务节点，并保留累计权重最高的 3 条路径。出边和路径使用

\begin{equation}
  W(p)=\sum_{e\in p} r_e\ln(1+c_e)
\end{equation}

排序，使高错误率且具有足够调用样本的下游边优先于偶发单次错误。到达异常子图叶节点或深度上限即停止。每次提取到的边还会写入 \filepath{scry\_memory\_topology\_edges}，累计观察次数、错误观察次数、平均错误率、置信度、最近 Trace 和时间；在线检索仍以本次证据中的运行边为准，历史拓扑不会伪装成本次事故中真实存在的调用关系。

对于路径 $p=(s_0,\ldots,s_n)$ 上的服务 $s_i$，拓扑得分定义为

\begin{equation}
  T(s_i,p)=35+65\frac{i}{n},\qquad n>0,
\end{equation}

因此越靠近异常路径终端的服务得到越高优先级，非路径服务得分为 0。该设计把“应该先查哪个服务”和“该服务可能是哪一种故障模式”拆成两个阶段，避免入口症状文档因词汇重合而长期压过根因服务记录。

\subsubsection{错误签名与稳定记忆分块}

错误签名由确定性规则提取，不调用模型。提取器读取问题文本及可用证据中的标题、摘要和结构化正文，识别 HTTP 状态码、gRPC/SQLSTATE 标识、以 error、exception、timeout、exhausted、refused、unavailable、overflow、deadlock 和 throttling 等模式结尾的错误词，以及“错误、异常、失败、超时、耗尽、拒绝、限流、熔断、延迟、重试”等中英文关键词。签名最多保留 24 项，既用于单路径终端支持校验，也用于多路径分支选择和候选分块重排。

每条记忆由摘要和结构化正文生成稳定分块。系统优先按段落和句子边界切分，单块不超过 1,800 字符，单条记忆最多 80 块；块类型自动归类为 summary、observation、procedure、evidence 或 metadata。每块保存 SHA-256、词项数、关键词、服务名和置信度，重复内容按哈希去除。长记忆因此可以在 UI 中逐块预览，检索也只把得分最高的 3 个块放入 Memory Packet，避免整篇文档挤占模型上下文。

对查询 $q$ 和记忆块 $b$，系统计算词项 Jaccard 相似度 $L(q,b)$、256 维哈希特征余弦相似度 $D(q,b)$ 和错误签名覆盖率 $G(b,\sigma)$。块内排序为

\begin{equation}
  B(q,b)=0.45L+0.35D+0.20G.
\end{equation}

该局部稠密通道完全在 Agent Service 内执行，不依赖外部 Embedding 服务；它作为中英文词形差异的补充，而不是替代运行拓扑。

\subsubsection{TopoMem-MP 与选择性回退}

当异常子图只有一条候选路径时，系统使用 TopoMem；存在多条路径时使用 TopoMem-MP。TopoMem-MP 分别统计每条路径终端服务记忆对错误签名的最大支持度：若某一路径获得唯一正支持，只保留该路径；若多个分支同分，保留并融合候选；若所有终端都没有签名支持，则放弃拓扑约束并进入词项回退。该机制用于处理重试、并发故障、旧错误分支和低流量噪声，避免简单贪心始终追随错误率最高的分支。

选择性回退还覆盖三种情况：当前证据没有可解析 Trace；存在边但无法从症状服务形成异常路径；单路径终端没有记忆，或终端记忆的签名/局部稠密支持低于 5 分。每次回退都返回明确 \code{fallbackReason}，不会静默把平面检索描述成拓扑定位。回退模式继续使用词项、局部稠密、签名和质量因子，使没有 Trace 的日志类事故仍能使用历史经验。

\subsubsection{多因子评分与可解释检索包}

候选质量分综合记忆置信度 $C$、重要性 $I$、活性 $V$ 和验证状态：

\begin{equation}
  Q=0.35C+0.25I+0.30V+10\mathbb{I}_{verified}.
\end{equation}

TopoMem 模式的最终分数和词项回退分数分别为

\begin{align}
  S_{topo}&=0.58T+0.17G+0.09L+0.07D+0.09Q,\\
  S_{fallback}&=0.46L+0.34D+0.10G+0.10Q.
\end{align}

结果先按最终分数、再按活性排序，默认返回 5 项，接口允许 1--30 项。每个结果同时返回 lexical、dense、topology、signature、quality 和 final 六个因子，匹配路径、命中签名、前三个记忆分块及其独立得分也一并返回。检索运行写入 \filepath{scry\_memory\_retrieval\_runs}，记录 strategy、symptomService、signature、paths、fallbackReason、候选数量、返回记忆 ID、耗时、Thread 和 Trace，使一次“为什么先查这个服务”的选择可以被复盘。

\subsubsection{验证写回、自学习与自维护}

TopoMem 仓库区分 working、episodic、semantic 和 procedural 四类记忆，以及 organization、user、agent、thread 和 workflow 五种作用域。组织级写入仅允许管理员；线程和 Loop 记忆必须绑定对应 ID。系统提供三条学习路径：用户或 Agent 通过 \code{remember} 把选定证据写回；证据被人工接受后自动生成 verified 经验；Loop 完成后把结果与可用证据共同写为 episodic 或 procedural 记忆。已驳回或挂起证据不能作为来源，来源表保存 evidence、message、workflow、manual、memory 或 document 标识、Trace、关系和权重。

写入前执行秘密遮蔽、内容稳定序列化、120,000 字节上限和 Prompt Injection 检测。相同内容哈希再次出现时不创建副本，而是增加强化次数、活性、置信度和重要性；新记忆会与同组织最近记录建立 duplicates、conflicts、supports 或 same-service 关系。用户可以固定、验证、挂起、归档或遗忘记忆；遗忘会清空正文、摘要、标签、关键词、服务标识和分块，仅保留不可混淆的审计骨架。

维护器为不同类型设置不同半衰期：working 为 1 天，episodic 为 90 天，semantic 为 365 天，procedural 为 540 天。未固定且未验证记忆的活性按

\begin{equation}
  V=0.34I+0.32C+0.24R+0.10F
\end{equation}

重新计算，其中 $R=100\cdot0.5^{d/h}$ 表示距最近访问或强化 $d$ 天后的时效分，$h$ 为类型半衰期，$F$ 由强化次数和访问次数共同生成。活性低于 20 的在线记忆自动归档；相似度不低于 0.88 的重复记忆合并并强化保留项；同主题出现明显相反状态时建立 conflicts 关系，并把较低置信度且未验证的记忆挂起等待人工复核。维护报告返回扫描、衰减、强化、合并、冲突、归档和过期数量，并写入安全审计。

记忆管理页面展示总量、在线量、验证量、挂起量、平均活性、类型和状态分布；详情视图展示正文、分块预览、关键词、来源、关系和完整维护历史。执行 TopoMem 检索后，页面进一步显示检索策略、路径、签名、六因子得分和命中块，因此记忆不是隐藏在 Prompt 中的黑盒上下文，而是可以逐项检查、删除、修订和重新验证的运行资产。

\subsection{MCP 工具协议与插件注册}

Agent 的实际工具调用链通过 Model Context Protocol 完成。内置 \code{/mcp} 端点采用 Streamable HTTP 传输，每个请求使用当前 Scry Bearer Token 重新解析用户、组织和角色，再创建请求作用域的 MCP Server。Server 当前发布 16 个工具，覆盖计划、服务、告警、仪表盘、Query Range、安全意图、SSH、证据链与 TopoMem。MCP Client 在创建 Agent 时执行 initialize 与 tools/list，将返回的 JSON Schema 转换为 Mastra 工具，并在对话流或 Loop 节点结束后主动断开连接。

运行时 tools/list 使用 camelCase 名称，例如 \code{querySignals}、\code{listEvidence} 与 \code{sshReadonlyInspect}；装载到 Agent 后由 MCP Client 增加服务命名空间，分别成为 \code{scry\_querySignals}、\code{scry\_listEvidence} 与 \code{scry\_sshReadonlyInspect}。系统指令与技术文档均使用这一运行时映射，避免模型依据内部源码 ID 猜测工具名。

每个 inputSchema 都是 \code{additionalProperties=false} 的对象，顶层和嵌套字段包含中文说明、枚举、默认值、边界、引用来源与调用示例。模型必须直接提交一个匹配 Schema 的 JSON 对象，不得再添加参数包装层；无参数工具固定提交空对象。ID 型字段必须来自前置工具真实输出，例如 SSH 主机和命令来自 \code{sshListHosts}，证据 UUID 来自 \code{listEvidence}，记忆 UUID 来自 \code{searchMemory}。校验失败后只修正被指出的字段，不允许切换为底层协议继续试探。考虑到系统支持自定义 OpenAI Responses 兼容端点，Scry 不发送 provider strict 扩展元数据，参数的确定性校验在 MCP Server 工具边界完成。

\code{querySignals} 使用单一扁平对象，唯一必填顶层字段 \code{signal} 取 \code{logs}、\code{traces} 或 \code{metrics}。公共字段包括回看分钟数、案例编号、服务列表、附加过滤、结果模式、分组和时间桶；日志增加正文搜索与级别，链路增加操作名、仅错误和最小耗时，指标增加指标名、指标类型、聚合和 temporality。模型不得添加外层 \code{query}，也不得提交 \code{compositeQuery}、\code{logQueries} 或 \code{builderQueries}。\code{queryType}、\code{start} 与 \code{end} 由编译器生成。

工具执行层把高层输入确定性编译为 Query Range V3：自动计算毫秒时间窗和 step，固定声明 \code{queryType=builder}，创建 \code{builderQueries.A}，并按记录或时间序列选择 panelType、聚合、排序、分页和字段作用域。日志与链路中的 \code{service.name} 映射到资源属性；Trace ID、Span ID、日志级别等别名映射到对应标准列；\code{scry.debug.case\_id} 也映射到资源属性；其他 OpenTelemetry 字段映射到普通属性。Zod 在网络调用前拒绝缺少 metricName、错误聚合、非数组 in/nin 和越界时间范围，避免模型以底层 400 响应进行协议试探。

内置 MCP Server 在 initialize 响应中发布统一调用协议，Agent 同时使用同一份受版本控制的协议常量构建系统提示。MCP Client 保持 Server Instructions 自动转发关闭，避免自定义 Responses 兼容端点因重复 Guidance 进入空推理循环，并防止第三方内容取得系统指令权限。管理员可以按组织注册远程 MCP 服务。注册记录包含名称、说明、Streamable HTTP URL、启用状态、超时和可选请求头；请求头使用 AES-256-GCM 加密，列表接口只返回 \code{hasHeaders}。Agent 不向外部插件转发 Scry JWT，插件只能获得管理员为该插件单独配置的认证头。远程服务连接失败时仅跳过该插件并记录错误，内置 Scry MCP Server 失败则终止 Agent 创建，避免在缺少核心安全工具时降级运行。

Loop 的 \code{mcp-tool} 节点直接引用命名空间化的 toolId 和结构化参数。\code{security-check}、\code{operation} 与 \code{verify} 节点分别承担强制策略解释、审批后受控执行和结果复核。系统初始化提供安全决策、磁盘爆满、僵尸进程、磁盘 I/O、配置漂移、服务恢复、网络暴露面、综合根因和 Prompt Injection 审查预设，其中系统安全 Loop 锁定，不能被修改或删除。

\subsection{运行时创建与系统指令}

\code{createScryAgent(user, settings, mode)} 每次请求创建用户相关 Agent，ID 为 \code{scry-operator-<userId>}。Memory 的 \code{lastMessages} 来自管理员配置，资源 ID 为用户 ID，线程 ID 为对话 ID，从而形成

\begin{equation}
  K_{memory} = (userId, threadId).
\end{equation}

该键只表示短期会话 Memory。TopoMem 是独立的长期运维记忆层，按 orgId、userId、scope、serviceName、status 和运行证据上下文执行检索，并把结果压缩为有界 Memory Packet。二者同时存在：短期 Memory 保持对话连续性，TopoMem 负责跨 Thread、跨 Loop 和跨事故复用已验证经验，证据引擎负责阻止两者被误当作当前事实。

系统指令固定要求使用简体中文、复杂诊断先调用 \code{scry\_publishPlan} 发布 2 至 6 步计划、仅依据工具数据下结论、查询必须有界、严格遵循 MCP inputSchema、不得猜测底层 HTTP API、SSH 只能使用已配置的主机和命令 ID、正文不重复全部 A2UI 原始数据，并采用“结论、证据、下一步”的输出顺序。工具校验失败、工具不可用、无查询结果和确认存在服务故障被定义为不同状态，模型不得把接口错误转换为故障事实。管理员补充指令附加在内置规则之后，不替换基础安全规则。

三种交互模式改变任务姿态：\code{diagnose} 优先用最少观测查询定位根因；\code{plan} 澄清目标、依赖、风险和执行步骤，不触发外部变更；\code{operate} 可以调用管理员启用的操作工具，高风险调用进入人工审批。模式只是行为约束，真正的执行权限仍由工具注册、用户角色、主机状态、命令策略和审批结果共同决定。

\subsection{Prompt 分层}

最终 Prompt 由六层内容组成。第一层是不可被管理员覆盖的 Scry 基础规则，包括语言、证据、权限和输出结构；第二层是模式规则；第三层是管理员补充指令；第四层是 TopoMem 检索包及其“必须重新验证”边界；第五层是 Thread 短期记忆和当前用户消息；第六层是本轮工具返回的观测证据。附件、记忆和工具结果都被标记为不可信数据，不被当作系统指令，其中只有第六层可以直接支撑本轮事实结论。

对于 Loop 中的 Agent 节点，编译器把节点 Prompt、用户输入和前序结果组合为节点级 Prompt。前序结果先序列化并按 50,000 字符截断，避免一个大型分支完全占据上下文。模型输出仅写入当前节点结果，后续节点通过 \code{results.<stepId>} 显式引用，减少隐式状态传播。

\subsection{模型供应商适配}

模型设置支持 \code{openai} 与 \code{anthropic} 两种协议，以及 \code{official} 和 \code{custom} 两类端点。OpenAI 路径固定使用 Responses API，Anthropic 路径使用 Messages API\cite{openairesponses,anthropicmessages}。自定义 URL 会去除末尾的 \code{/responses} 或 \code{/messages} 后再由适配器拼接，避免重复路径。

每次生成应用温度、最大输出 Token、最大步骤、失败重试和工具选择策略。\code{toolChoice=auto} 由模型判断是否调用工具，\code{none} 禁止工具，\code{required} 要求至少一次工具调用。\code{maxSteps} 同时限制模型与工具的往返次数，\code{maxRetries} 只处理供应商级可重试错误，不重复已产生副作用的工具。

模型连接测试使用独立 AbortController 和超时，要求模型只回复“连接成功”，并将失败归类为配置、认证、模型、限流、超时、网络或供应商问题。错误返回前隐藏疑似 API Key 和外部 URL，并截断为 500 字符。连接测试使用当前表单值和已保存密钥的合并结果，因此管理员可以在保存前验证新模型、端点和凭据。

模型配置是运行时读取的实例级设置。新的请求使用最新配置，已经开始的 Run 保持创建时的模型实例和参数，避免中途切换供应商改变同一执行的语义。Endpoint Preview 在前端展示最终请求 URL，便于排查根路径、\code{/responses} 或 \code{/messages} 重复拼接问题。

\subsection{对话、附件与流式持久化}

对话 API 的执行序列如图~\ref{fig:chat-sequence}。线程所有权同时校验 \code{threadId}、\code{userId} 和 \code{orgId}。用户消息和可解析文本附件先进入安全决策链，拒绝结果返回 traceId 且不进入模型；允许或需审批的请求再写入 UI 消息表，创建 MCP 工具运行时并调用 \code{handleChatStream}。流被 \code{tee()} 分为客户端分支和持久化分支，最终 Assistant UIMessage 以 Upsert 方式写入 WASM SQLite。

\begin{figure*}[t]
\centering
\begin{tikzpicture}[
  actor/.style={draw=scryline,fill=scrygray,minimum width=24mm,minimum height=7mm,align=center,font=\small},
  msg/.style={-{Latex[length=1.7mm]},draw=scryline},
  lifeline/.style={draw=scryline,dashed},
  every node/.style={font=\scriptsize}
]
  \node[actor] (u) at (0,0) {Browser};
  \node[actor] (a) at (3.7,0) {Agent API};
  \node[actor] (q) at (7.4,0) {Query Service};
  \node[actor] (m) at (11.1,0) {Model/Mastra};
  \node[actor] (d) at (14.8,0) {SQLite / PostgreSQL};
  \foreach \x in {u,a,q,m,d} {\draw[lifeline] (\x.south) -- ++(0,-5.5);}
  \draw[msg] (u.south) ++(0,-0.5) -- node[above]{POST /api/chat + JWT} (a.south |- 0,-1.2);
  \draw[msg] (a.south) ++(0,-1.1) -- node[above]{JWT 校验与角色查询} (q.south |- 0,-1.8);
  \draw[msg] (q.south) ++(0,-1.7) -- node[above]{角色/有效性} (a.south |- 0,-2.4);
  \draw[msg] (a.south) ++(0,-2.3) -- node[above]{保存用户消息} (d.south |- 0,-3.0);
  \draw[msg] (a.south) ++(0,-2.9) -- node[above]{messages + memory key} (m.south |- 0,-3.6);
  \draw[msg] (m.south) ++(0,-3.5) -- node[above]{tool call + 当前 JWT} (q.south |- 0,-4.2);
  \draw[msg] (q.south) ++(0,-4.1) -- node[above]{有界观测结果} (m.south |- 0,-4.8);
  \draw[msg] (m.south) ++(0,-4.7) -- node[above]{UI Message Stream} (u.south |- 0,-5.4);
  \draw[msg] (a.south) ++(0,-4.8) -- node[above]{异步保存最终消息} (d.south |- 0,-5.4);
\end{tikzpicture}
\caption{Agent 对话请求、授权继承、工具取证和流式持久化时序。}
\label{fig:chat-sequence}
\end{figure*}

每条消息最多 5 个附件，单文件不超过 10 MB。\code{txt/md/json/log/csv} 等文本附件不超过 512 KB，并在送入模型前转换为带“用户提供的数据”边界标记的文本。文件名会移除换行和方括号，避免破坏边界结构；其他文件保留为模型文件部件，由供应商能力决定处理方式。Loop 输入只接受结构化文本和 JSON，文件型任务先在普通 Thread 中解析，再把结果传入 Loop。

\subsection{UI Message Stream 协议}

Scry 不把模型响应简化为单一文本，而是完整传输 AI SDK UI Message Part。常见 Part 包括 \code{text}、\code{reasoning}、\code{source-url}、\code{source-document}、\code{file}、\code{step-start}、\code{tool-*}、\code{data-tool-call-approval}、\code{data-scry-workflow-step}、\code{data-scry-workflow-result} 和 \code{data-tripwire}。每个 Part 都有稳定 type，工具 Part 还包含 state、input、output 和 errorText。

前端依据 Part state 显示输入流式、等待审批、执行中、输出可用、失败、暂停和取消。最终 UI Message 以完整 JSON 保存，因此页面刷新后 reasoning、来源、工具输入输出和审批卡片仍然存在。重新生成请求携带 trigger 和 messageId，后端不会重复保存同一用户消息；审批恢复请求携带 runId 与 resumeData，继续原 Run 而不是创建新 Run。

\subsection{Scry 观测工具与 A2UI}

工具使用 Mastra \code{createTool} 与 Zod 同时声明输入和输出。输入 Schema 约束模型能够提交的字段，执行器依据认证上下文和工具策略完成二次校验，输出 Schema 则保证模型与前端接收稳定数据结构。表~\ref{tab:tools} 列出 Agent 运行时注册的标准工具。

\begin{table*}[t]
\caption{Agent 工具目录与控制措施}
\label{tab:tools}
\centering
\small
\begin{tabularx}{\textwidth}{p{2.5cm}p{2.7cm}Y p{2.5cm}}
\toprule
\textbf{工具 ID} & \textbf{后端操作} & \textbf{用途与输出} & \textbf{控制措施} \\
\midrule
\code{publishPlan} & 无外部查询 & 发布诊断步骤，生成 plan 型 A2UI & 1--8 步；字段完整描述；拒绝额外字段 \\
\code{listServices} & GET \code{/api/v1/services/list} & 服务清单和 table 型 A2UI & 继承 JWT/RBAC；只读；空对象输入 \\
\code{listAlerts} & GET \code{/api/v1/alerts} & 当前告警和 table 型 A2UI & 继承 JWT/RBAC；只读；空对象输入 \\
\code{listDashboards} & GET \code{/api/v1/dashboards} & 仪表盘清单和 table 型 A2UI & 继承 JWT/RBAC；只读；空对象输入 \\
\code{querySignals} & POST \code{/api/v3/query\_range} & 扁平高层格式查询指标、日志或链路，生成 key-value A2UI & 顶层 signal；Schema 预校验；自动编译 V3；200 KB 上限 \\
\code{evaluateSecurity}\newline\code{Intent} & 本地策略引擎 & 返回意图、风险分和 allow/approval/deny & text 与 mode 明确；不依赖外部模型 \\
\code{listEvidence} & 读取证据仓库 & 返回置信度、关联度、状态、选择和备注 & 默认排除 rejected/suspended \\
\code{getEvidenceChain} & 读取证据关系 & 返回证据节点与印证、派生关系 & 对话与组织双重隔离 \\
\code{searchMemory} & TopoMem 检索 & 返回根服务优先的候选记忆与评分 & query 可省略；limit 为 1--20 \\
\code{remember} & 写入记忆仓库 & 写入证据支撑的可复用运行经验 & 正文为字符串；证据 UUID 必填；秘密检测 \\
\code{inspectMemory} & 读取记忆仓库 & 返回正文、分块、来源、关系和维护历史 & memoryId 来自检索结果 \\
\code{maintainMemory} & 记忆维护器 & 执行衰减、强化、去重、冲突和归档 & 批量上限 500；不硬删除 \\
\code{forgetMemory} & 记忆遗忘器 & 清空可检索内容并保留审计骨架 & 精确确认字符串 FORGET \\
\code{sshListHosts} & 读取 Agent Repository & 返回启用主机、只读命令 ID、扩展命令元数据 & 不返回密钥或精确命令文本 \\
\code{sshReadonly}\newline\code{Inspect} & 执行内置固定命令 & 系统、磁盘、内存、失败服务、监听端口、近期错误 & hostId/commandId 来自目录；审批策略；审计 \\
\code{sshExecute}\newline\code{Command} & 执行管理员预存命令 & 执行精确命令定义，返回 code 型 A2UI & 拒绝主机名和原始 Shell；逐次审批；审计 \\
\bottomrule
\end{tabularx}
\end{table*}

\code{callScry} 始终把当前用户 JWT 放入 Authorization Header，并在网络层设置连接、响应和总体超时。Query Service 返回非 2xx 状态时，工具将响应转换为统一错误对象；响应体超过配置上限时，执行器保留元数据、字段结构、前后样本和截断标记。工具输出同时包含原始 \code{data}、展示结构 \code{a2ui}、关联标识和执行统计，模型可以依据原始数据推理，用户则获得适合扫描、比较和审批的视觉结果。

\subsection{工具调用状态机与 A2UI Schema}

每个工具调用使用 \code{toolCallId} 标识，并依次经历 \code{input-streaming}、\code{input-available}、\code{approval-requested}、\code{output-available} 或 \code{output-error} 状态。只读且无需审批的工具从输入可用直接进入执行；受控工具先生成审批部件，批准后以相同 \code{toolCallId} 恢复。取消、超时和策略拒绝都产生终态事件，前端不会把它们显示为仍在运行。

A2UI 0.1 顶层包含 \code{version}、\code{surfaceId}、\code{operation} 和 \code{root}。其中 \code{surfaceId} 在 Thread 内唯一，用于替换或增量更新同一结果。\code{operation} 支持 replace、append 和 update；根节点类型及其语义见表~\ref{tab:a2ui-types}。所有类型共用标题、数据、敏感字段标记、截断标记、来源和时间范围元数据。

\begin{table}[t]
\caption{A2UI 组件语义}
\label{tab:a2ui-types}
\centering
\small
\begin{tabularx}{\columnwidth}{p{1.9cm}Y}
\toprule
\textbf{类型} & \textbf{渲染与使用方式} \\
\midrule
\code{plan} & 展示步骤序号、名称、执行状态和当前步骤，用于诊断计划与 Loop 进度。 \\
\code{table} & 展示列定义、行数据、单位、排序和分页摘要；大结果提供截断说明。 \\
\code{key\_value} & 展示指标、属性、查询摘要和主机巡检结果，支持单位格式化。 \\
\code{status} & 展示成功、警告、失败、暂停、取消及其说明。 \\
\code{code} & 展示预配置命令 ID、退出码和受控 stdout/stderr 片段。 \\
\code{markdown} & 展示结构化说明、结论和引用，不执行嵌入式脚本。 \\
\bottomrule
\end{tabularx}
\end{table}

前端渲染前再次验证 A2UI 版本和根节点类型，所有文本按纯内容处理，Markdown 禁止原始 HTML，表格单页按显示上限裁剪。无法识别的类型进入通用 JSON Inspector，并保留原始 Part，保证协议升级期间的数据兼容性。A2UI 与模型正文相互独立：A2UI 用于证据和操作状态，正文用于解释和结论，避免模型复制大段原始数据造成上下文与页面冗余。

\subsection{SSH 执行链路}

SSH 纳管由 \code{scry-ops} 环境安装、主机登记、指纹核验、私钥加密、策略绑定、连接测试和启用七个步骤组成。受管主机固定使用 \code{scry-ops} 公钥账户，不接受 root、UID 0、其他用户名或密码认证。管理员先调用 probe 端点读取服务器公开主机密钥指纹，再通过可信运维通道核验，保存后连接器使用 \code{hostVerifier} 对 SHA-256 指纹做常量时间比较。指纹变化会阻断连接并产生安全审计事件，必须由管理员重新核验和更新。

私钥和 passphrase 以 AES-256-GCM 密文保存在 WASM SQLite。运行时只在建立连接前解密到进程内存，连接关闭后释放引用；列表、详情、审计和错误响应都不返回秘密。连接测试也通过固定的 \code{scry-connection-test} 包装器和 sshd ForceCommand，不存在测试旁路。命令执行禁用 PTY、交互输入、X11、端口转发和隧道，默认 30 秒超时，标准输出与标准错误合计捕获 200 KB，同时记录真实字节数和截断状态。

SSH 策略分为只读巡检和扩展执行。只读巡检提供七个内置命令 ID，全部映射到 \filepath{/opt/scry-ops/bin} 下由 root 持有的固定包装器：

\begin{itemize}
  \item \code{system\_overview} 与 \code{disk\_usage}；
  \item \code{memory\_usage} 与 \code{failed\_services}；
  \item \code{network\_listeners} 与 \code{recent\_errors}。
  \item \code{security\_preflight} 校验账户、UID、附加组、包装器属主和权限。
\end{itemize}

远端 sshd 对 \code{scry-ops} 使用 \code{ForceCommand /opt/scry-ops/bin/scry-dispatch}，并禁止 TTY、转发、隧道和用户环境。分发器只接受固定只读包装器或 \code{sudo -n} 调用的变更包装器；sudoers 仅允许精确包装器路径，服务名和清理时间等参数在分发器和包装器内各校验一次。扩展命令保存与执行前都经过词法分析，Shell 解释器、管道、重定向、命令替换、提权程序、破坏性程序和非包装器路径一律拒绝。

扩展命令固定进入人工审批。审批卡片包含主机名称、命令名称、风险等级、执行原因、超时和请求用户，不展示凭据。批准结果绑定 \code{runId}、\code{toolCallId}、\code{hostId}、\code{commandId} 和策略版本，任何字段变化都会使原批准失效。执行结束后，系统写入用户、组织、线程、运行、主机、工具、命令、策略版本、状态、退出码、耗时、输出字节数、截断标记和错误摘要；审计事件与 Agent Run 的 Trace 标识关联，可从对话、运行历史和安全审计三处定位同一次操作。

\begin{figure}[t]
\centering
\begin{tikzpicture}[
  node distance=3.2mm,
  every node/.style={font=\scriptsize},
  state/.style={draw=scryline,rounded corners=2pt,minimum width=0.76\columnwidth,minimum height=5.5mm,align=center,fill=white},
  arrow/.style={-{Latex[length=1.5mm]},draw=scryline}
]
  \node[state,fill=blue!7] (request) {工具请求：hostId + commandId};
  \node[state,below=of request] (policy) {角色、主机、指纹、策略和命令版本校验};
  \node[state,fill=orange!8,below=of policy] (approval) {生成审批并暂停 Run};
  \node[state,below=of approval] (execute) {批准后建立 SSH 连接并执行固定命令};
  \node[state,fill=green!8,below=of execute] (audit) {返回 A2UI，写入审计与 OpenTelemetry Span};
  \draw[arrow] (request) -- (policy);
  \draw[arrow] (policy) -- (approval);
  \draw[arrow] (approval) -- (execute);
  \draw[arrow] (execute) -- (audit);
\end{tikzpicture}
\caption{SSH 受控执行闭环。任一校验失败都会终止调用并写入拒绝事件。}
\label{fig:ssh-execution}
\end{figure}

\subsection{Loop 定义语言}

Loop 是持久化、版本化且可恢复的工作流定义。定义包含基本信息、输入 Schema、触发器和步骤树；步骤树使用版本 1/2 DSL 表示。叶节点负责执行具体能力，控制节点负责组织分支与迭代。

\begin{itemize}
  \item \textbf{叶节点：}推理节点为 \code{agent}，MCP 调用节点为 \code{mcp-tool}，强制安全节点为 \code{security-check}，受控执行和结果复核节点为 \code{operation}/\code{verify}，人工控制节点为 \code{approval}，观测读取节点为 \code{services}、\code{alerts}、\code{dashboards}。
  \item \textbf{条件节点：}\code{condition} 使用路径、操作符和值，分别编译 true/false 子流程。
  \item \textbf{并行节点：}\code{parallel} 将 2 至 6 个分支编译为 Mastra \code{parallel}，完成后合并上下文。
  \item \textbf{循环节点：}\code{loop} 支持 \code{while}、\code{until} 和 \code{foreach}；foreach 可指定源路径和 1--10 的并发数。
\end{itemize}

条件路径从工作流上下文读取，操作符支持 \code{exists}、\code{equals}、\code{not\_equals}、\code{gt/gte/lt/lte}。路径以 \code{input}、\code{trigger}、\code{results}、\code{item} 或 \code{runtime} 为根，比较值按布尔、数字或字符串语义处理。工作流整体最多 100 个节点，控制结构最多嵌套 5 层，所有节点 ID 必须唯一；循环或 foreach 最大 50 次或 50 项，并行分支为 2 至 6 个，foreach 并发为 1 至 10。

Loop 输入字段支持 text、number、boolean 和 json。启动前，运行入口将用户输入、默认值与触发器载荷合并，再按 required 和字段类型校验。运行上下文采用如下结构：

\begin{equation}
 C_{loop}=\{input,trigger,results,item,runtime\}.
\end{equation}

其中 \code{input} 在 Run 生命周期内只读，\code{trigger} 保存 manual/cron/event 来源与元数据，\code{results} 按 stepId 保存步骤输出，\code{item} 只在 foreach 子流程中存在，\code{runtime} 保存 runId、workflowId、workflowVersion、principal、attempt 和计划时间。显式命名空间使条件表达式和后续 Prompt 不依赖隐含变量。

\subsection{Loop 编译与执行}

编译器递归执行 \code{compileSequence(record, definitions, principal, settings, id)}。保存 Loop 时先执行结构验证，启动 Run 时再按固定版本编译，保证运行中的版本不受后续编辑影响。普通节点通过 \code{then} 串联；条件节点通过 \code{branch}；并行节点通过 \code{parallel}；while 和 until 分别映射到 \code{dowhile} 与 \code{dountil}；foreach 先读取源数组、应用最大项数，再映射到有界并发子流程。每个控制节点前后插入标记步骤，用于发布进度、记录分支选择并把输出归一化为统一上下文。

\begin{figure}[t]
\centering
\begin{tikzpicture}[
  node distance=4mm,
  every node/.style={font=\scriptsize},
  state/.style={draw=scryline,rounded corners=2pt,minimum width=23mm,minimum height=6mm,align=center,fill=white},
  arrow/.style={-{Latex[length=1.6mm]},draw=scryline}
]
  \node[state] (queued) at (0,0) {queued\\等待执行};
  \node[state] (running) at (3,0) {running\\步骤执行};
  \node[state,fill=orange!8] (suspended) at (6,0) {suspended\\等待审批};
  \node[state,fill=green!8] (completed) at (6,-1.4) {completed\\写入输出};
  \node[state,fill=red!8] (failed) at (3,-1.4) {failed/cancelled\\终态记录};
  \draw[arrow] (queued) -- (running);
  \draw[arrow] (running) -- node[above]{approval} (suspended);
  \draw[arrow] (suspended) -- node[right]{approve} (completed);
  \draw[arrow] (suspended) -- node[below]{reject} (failed);
  \draw[arrow] (running) -- node[below]{success} (completed);
  \draw[arrow] (running) -- (failed);
\end{tikzpicture}
\caption{Loop 运行状态与人工审批恢复。}
\label{fig:workflow-state}
\end{figure}

叶节点执行前发布 \code{data-scry-workflow-step/running}，并把 attempt、startedAt 和输入摘要写入步骤状态。\code{agent} 节点将节点 Prompt、结构化输入和前序 \code{results} 的有界摘要组合为请求；观测节点调用同一个 Scry Read Operation；审批节点调用 Mastra \code{suspend}，保存 resume label、审批快照和到期时间。步骤成功后写入 output、duration 和 completedAt；失败时依据节点重试次数分类为可重试或终态错误，并使用指数退避重新入队。

并行分支在独立子上下文中执行，完成后按 branchId 合并；任一分支失败时，控制节点按照 all-success 语义结束并记录各分支状态。while/until 每次迭代后重新计算条件，foreach 为每个 item 生成稳定 itemKey，结果按原数组顺序归并。最终步骤发布 \code{data-scry-workflow-result}，写入输出摘要、完整结构化输出和运行统计；与 Thread 关联的 Run 还会把最终文本追加到对话历史。

\subsection{Cron 调度器}

Schedule Coordinator 在 Agent Service 启动后加载所有已启用且包含 Cron 的 Loop。Cron 表达式在保存时完成语法校验，时区使用 IANA 标识；调度器将本地日历时间转换为 UTC 计划时间，并在夏令时切换时按时区规则确定唯一触发点。每次 Loop 变更都会刷新内存注册表，停用或删除操作同步撤销任务。

多副本部署使用租约选主。协调器周期性竞争 \code{scheduler-leader} 租约，租约包含 holderId、epoch 和 expiresAt；只有持有当前 epoch 的副本可以扫描到期任务。扫描器为每个计划时间生成

\begin{equation}
 K_{cron}=H(w,v,t_s),
\end{equation}

其中 $w$ 为 workflowId，$v$ 为 workflowVersion，$t_s$ 为 scheduledAt。

Trigger Claim 表对 \code{triggerKey} 设置唯一约束。领取成功的任务创建 \code{queued} Run，领取冲突表示其他副本已经处理，因此直接跳过。服务停机跨过触发点时，恢复扫描按照 misfire policy 处理：\code{skip} 跳过过期点，\code{fire-once} 补执行最近一次，\code{catch-up} 在最大补偿次数内逐次创建 Run。Scry 默认采用 fire-once，避免长时间停机后形成执行洪峰。

调度身份在启用 Loop 时固化为 \code{executionPrincipal}，包含组织、所有者和权限副本。触发时 Coordinator 向 Query Service 的内部身份端点换取短期 JWT，工具继续通过 Query Service RBAC 访问数据。所有 Cron Run 都记录 scheduledAt、claimedAt、startedAt、principal、schedulerInstance 和 triggerKey，便于定位延迟与重复抑制结果。

\subsection{事件触发器}

Event Coordinator 接收平台内部事件并转换为规范事件信封。事件信封包含 eventId、eventType、occurredAt、orgId、source、subject、attributes、payload 和 traceContext。Query Service、告警管理器及其他集成通过带服务身份的 \code{/api/internal/events} 提交事件；接收端验证签名、时间戳和重放窗口后，先写入 Event Inbox，再异步匹配 Loop，保证调用方无需等待模型执行。

每个事件触发器由 eventType、enabled 和 filter 组成。eventType 支持精确名称和受控前缀通配；filter 被编译为只读表达式，只能访问信封字段，支持 equals、not-equals、contains、in、gt/gte/lt/lte 和 exists。保存触发器时完成语法与字段校验，运行时不执行动态代码。匹配成功后生成

\begin{equation}
 K_{event}=H(workflowId,triggerId,eventId).
\end{equation}

唯一 Trigger Claim 保证同一事件对同一触发器只创建一次 Run。事件载荷写入 \code{trigger.payload}，同时应用大小上限、敏感字段遮蔽和类型校验。Inbox 消费成功后记录 matched、ignored 或 failed；暂时性失败按退避策略重试，达到次数上限后进入 dead-letter 状态并保留错误分类。管理员可在运行历史中按 eventId 重新投递，重新投递生成新的 replayId，同时保留原始 eventId 关联。

\subsection{并发控制、幂等与重试}

Run Dispatcher 在创建 Run 前应用三层并发控制：实例级总并发、组织级并发和 Loop 级并发。超过配额的 Run 保持 queued，并按优先级、计划时间和创建时间调度。Loop 支持 \code{allow}、\code{forbid} 和 \code{replace} 三种并发策略：allow 允许多个 Run 并行；forbid 在已有活动 Run 时保留一个待执行实例；replace 取消旧 Run 并启动最新实例。

幂等覆盖触发和步骤两个层级。触发层以 Trigger Claim 防止重复创建 Run；步骤层为具有外部副作用的调用生成

\begin{equation}
 K_{step}=H(r,s,a,c),
\end{equation}

其中 $r$、$s$、$a$、$c$ 分别表示 runId、stepId、attempt 和 toolCallId。Repository 先登记 executing，再执行工具，完成后保存输出摘要。供应商网络重试只覆盖尚未产生工具输出的模型请求；SSH 和其他副作用工具不做透明重放，必须依据已保存执行状态人工确认。

错误分为 validation、authentication、authorization、rate-limit、timeout、dependency、provider、tool、approval-rejected、cancelled 和 internal。validation/authorization/approval-rejected 直接结束；rate-limit、timeout、dependency 和 provider 按配置重试；internal 记录 traceId 并由运行器判定。退避时间采用带抖动的指数函数

\begin{equation}
 d_n=\min(d_{max},d_0 2^n)+U(0,j),
\end{equation}

其中 $n$ 为重试序号，$j$ 为抖动窗口。所有重试沿用同一个 runId，attempt 单调增加，便于用户区分一次业务运行和多次执行尝试。

\subsection{审批暂停与恢复}

审批节点进入 suspended 状态时，运行器提交当前步骤状态、上下文摘要、Workflow Version、resume label 和审批载荷，随后释放计算资源。审批记录包含 approver scope、requestedAt、expiresAt、decision、comment 和 decidedAt。审批 API 校验运行所有权、审批角色、当前 suspendedStep、版本和状态，使用条件更新保证同一审批只能成功决策一次。

批准后，\code{/api/workflow-runs/:runId/resume} 将 resumeData 写入恢复事件，Run Dispatcher 重新装载固定版本定义和持久化上下文，从相应 resume label 继续；拒绝后运行进入 failed，错误分类为 approval-rejected；超时后进入 expired。恢复不会重新执行审批点之前已经完成的步骤。服务重启时 Recovery Worker 扫描 running、retrying 和 suspended Run：running 根据最后一个持久化检查点重新排队，retrying 根据 nextAttemptAt 恢复，suspended 保持等待决策。

\subsection{Agent HTTP API}

\begin{table*}[t]
\caption{Agent Service API 分组}
\label{tab:agent-api}
\centering
\small
\begin{tabularx}{\textwidth}{p{2.1cm}p{5.5cm}Y}
\toprule
\textbf{分组} & \textbf{端点} & \textbf{授权与语义} \\
\midrule
健康 & GET \code{/health} & 唯一免认证端点 \\
模型 & GET/PUT \code{/api/settings/model}；POST \code{/test} & 读取隐藏密钥；更新和测试仅 ADMIN \\
SSH 主机 & GET/POST \code{/api/ssh/hosts}；probe、test、PUT、DELETE & 全部仅 ADMIN；指纹探测、连接测试均有显式 API \\
SSH 策略/审计 & GET/PUT \code{/api/ssh/policies}；GET \code{/api/ssh/audits} & 仅 ADMIN；审计查询上限 200 条 \\
线程 & GET/POST \code{/api/threads}；GET/DELETE \code{/:threadId} & 按用户和组织隔离，删除级联 UI 消息与关联运行 \\
对话 & POST \code{/api/chat} & 线程所有权校验；诊断/规划/操作；流式输出和恢复 \\
证据 & GET \code{/api/threads/:threadId/evidence}；PATCH \code{/api/evidence/:evidenceId}；POST \code{/evidence/revise} & 自动回填、状态/选择/备注更新和基于选定证据重新回答 \\
TopoMem & GET/POST/PATCH/DELETE \code{/api/memories}；search、stats、maintenance、samples & 作用域隔离、检索解释、维护、遗忘及论文语料示例生成 \\
Loop 定义 & GET/POST/PUT/DELETE \code{/api/workflows} & 查询登录可用；写操作仅 ADMIN \\
Loop 运行 & GET \code{/runs}；GET \code{/runs/:runId}；POST \code{/stream}；POST \code{/workflow-runs/:runId/resume|cancel} & 按用户和组织过滤；返回流、步骤状态、触发信息和恢复语义 \\
调度管理 & GET \code{/api/scheduler/status}；POST \code{/api/workflows/:id/trigger} & ADMIN 查看租约、队列和下次执行时间；手动触发复用统一 Run 管线 \\
内部事件 & POST \code{/api/internal/events}；POST \code{/api/internal/events/:eventId/replay} & 服务身份认证；写入 Inbox、匹配触发器、幂等创建 Run \\
运行事件 & GET \code{/api/workflow-runs/:runId/events} & 返回状态转换、步骤、工具、审批、重试和调度事件 \\
\bottomrule
\end{tabularx}
\end{table*}

所有 JSON API 返回 \code{requestId}。成功响应使用 \code{data} 与可选 \code{meta}；失败响应使用 \code{error.code}、\code{error.message}、\code{error.details} 和 \code{traceId}。流式接口在 HTTP 建立后通过结构化 Part 报告业务错误，客户端取消通过 Abort Signal 传播到模型和只读查询。创建操作返回 201，幂等删除返回 204，参数错误返回 400，身份问题返回 401，权限问题返回 403，对象不存在返回 404，版本或状态冲突返回 409，限流返回 429，依赖超时返回 504。

\subsection{Agent 持久化模型}

Agent Service 的持久化模型按八个领域组织：
\begin{itemize}
  \item 模型配置：\filepath{scry_model_settings}；
  \item 对话：\filepath{scry_agent_threads}、\filepath{scry_agent_ui_messages}；
  \item 证据：\filepath{scry_evidence_items}、\filepath{scry_evidence_relations}、\filepath{scry_evidence_actions}、\filepath{scry_evidence_revisions}；
  \item TopoMem：\filepath{scry_memory_items}、\filepath{scry_memory_chunks}、\filepath{scry_memory_sources}、\filepath{scry_memory_relations}、\filepath{scry_memory_actions}、\filepath{scry_memory_topology_edges}、\filepath{scry_memory_retrieval_runs}；
  \item SSH：\filepath{scry_ssh_hosts}、\filepath{scry_tool_policies}、\filepath{scry_ssh_audit_logs}；
  \item Loop：\filepath{scry_workflows}、\filepath{scry_workflow_runs}、\filepath{scry_workflow_run_events}；
  \item 触发：\filepath{scry_trigger_claims}、\filepath{scry_event_inbox}；
  \item 协调：\filepath{scry_runtime_leases}、\filepath{scry_agent_audit_events}。
\end{itemize}
模型设置以组织和环境为作用域，密钥只保存密文与 keyVersion。Thread、Message、Run、Workflow、SSH Host 和策略都包含 orgId；用户访问采用 orgId 与 userId 双重过滤，平台管理员使用显式管理端点跨用户查询。Workflow 保存当前定义和版本索引，Run 固定引用 workflowVersion；Trigger Claim 以 triggerKey 唯一约束实现幂等；Runtime Lease 使用 name 与 epoch 保证选主正确性。

UI Message 作为完整 JSON 保存，以适配 AI SDK Part 的扩展。工作流定义、输入、输出、步骤状态、触发器和输入模式也使用 JSON 列，但关键检索字段保持独立列，包括 orgId、userId、workflowId、runId、status、triggerType、scheduledAt、nextAttemptAt、createdAt 和 updatedAt。写入路径同时经过 Zod 校验、数据库约束和状态转换检查；Run Event 采用只追加模型，保存每次状态变化的序号、事件类型、时间、主体、摘要和 traceId。

Repository 使用短事务更新 Run 状态和追加事件。状态更新带 expectedStatus 条件，防止两个工作进程同时推进同一 Run；Trigger Claim 和租约更新依赖唯一键与 compare-and-swap 语义。大工具结果不重复写入事件表，事件只保存摘要和内容哈希，完整结果保存在 UI Message 或 Run Output 中。

\section{安全架构}

\subsection{身份认证与请求验证}

浏览器通过 Scry 登录流程获得 Bearer JWT，JWT 声明遵循 RFC 7519\cite{rfc7519}。Query Service 负责签名、有效期、颁发方、受众和用户状态验证；Agent Service 在本地完成令牌格式检查后，使用同一令牌向 Query Service 执行在线身份确认。除 \code{/health} 外，所有 Agent API 都要求有效身份；Query Service 不可达、令牌过期、用户停用或组织不匹配时一律拒绝请求。

内部调用使用独立服务身份。Agent 向 Query Service 发起用户数据查询时透传当前用户 JWT；Cron 和事件 Run 通过 executionPrincipal 换取短期组织级 JWT；Query Service、Alertmanager 和集成组件调用事件入口时使用服务令牌和请求签名。用户令牌与服务令牌具有不同 audience，不能互换。每次访问都重新验证主体、资源和动作，符合零信任架构的持续验证原则\cite{nistzt}。

\subsection{授权与租户隔离}

平台角色为 ADMIN、EDITOR 和 VIEWER。VIEWER 可以查询遥测、查看仪表盘并运行获得授权的 Loop；EDITOR 可以创建和修改观测对象；ADMIN 可以管理用户、组织、模型、SSH、工具策略、Loop、调度器和审计。高风险操作还需满足工具级策略和人工审批，角色不会绕过审批。

所有业务对象带 orgId。Thread 和 Run 同时带 userId，用于个人任务隔离；Workflow、模型设置、SSH 主机和策略按组织管理；系统级管理员通过独立管理端点执行跨组织运维。Repository 查询把 orgId 作为强制条件，API 不接受客户端覆盖认证上下文中的组织。对象 ID 使用不可预测 UUID，所有权校验仍然在读取和写入时执行，避免把随机 ID 当作授权机制。

\subsection{Prompt 与工具安全边界}

Scry 将系统规则、模式规则、管理员指令、用户输入、附件、观测数据和工具输出放入不同语义边界。附件和遥测内容使用明确的数据标记，并在进入 Prompt 前应用大小、类型和敏感字段策略；任何出现在数据中的“指令”都不改变系统层规则。模型输出只表达调用意图，实际执行由服务端 Schema、RBAC、工具注册表、资源所有权和策略引擎决定。这一设计直接覆盖 Prompt Injection、过度代理和不安全工具调用等 LLM 应用风险\cite{owaspllm}。

输入安全引擎先执行 NFKC 归一化和零宽字符清理，再检测指令覆盖、系统提示词或凭证提取、安全与审批绕过、授权角色混淆、不可信工具输出提升为指令以及编码执行载荷。检测器输出 finding 列表和 0--100 风险分；意图分类器把请求归入 observe、diagnose、change、privilege-escalation、secret-access 或 unknown。策略引擎将两类结果合并为 \code{allow}、\code{require-approval} 或 \code{deny}。提权、凭证读取和高风险注入直接拒绝；变更请求进入审批；只读诊断继续执行。

工具侧采用第二道独立过滤。命令分析器解析引号、转义和 Token，拒绝分号、管道、重定向、命令替换、换行、Shell 解释器、\code{sudo su}、网络反弹程序和破坏性磁盘/账户操作；只有固定 \code{scry-ops} 包装器或 \code{sudo -n} 包装器形式可以进入审批。执行前策略同时校验主机用户名、toolId、commandId、命令哈希和策略结果，因此模型或自定义 Loop 不能绕过底层强制层。

统一审计表 \filepath{scry\_security\_audit\_events} 使用 traceId 串联以下事件组：
\begin{itemize}
  \item 输入与分类：\path{intent.received}、\path{prompt_injection.scanned}、\path{intent.classified}；
  \item 策略与审批：\path{policy.*}、\path{approval.*}；
  \item 工具执行：\path{command.risk_evaluated}、\path{tool.execution.*}。
\end{itemize}
事件只保存输入 SHA-256、结构化原因、风险分、目标和耗时，不保存 API Key、私钥、MCP 认证头或完整 Prompt。管理界面以时间线展示同一 Trace 的完整决策与执行链。

工具注册表为每个工具保存 riskLevel、requiredRole、approvalMode、timeout、resultLimit、allowedResources 和 policyVersion。调用前，Policy Evaluator 使用不可变请求上下文计算 allow、deny 或 require-approval。查询工具只接受结构化查询对象；SSH 工具只接受主机与命令 ID；模型不能传入数据库连接、文件路径、外部 URL 或原始 Shell。操作模式只改变可候选工具集合，不改变底层授权。

\subsection{秘密、密钥与网络控制}

模型 API Key 和 SSH Secret 通过 AES-256-GCM 加密。每次加密生成 12 字节随机 IV，密文格式为

\begin{equation}
  B = v \parallel keyId \parallel \mathrm{b64url}(IV) \parallel
  \mathrm{b64url}(Tag) \parallel \mathrm{b64url}(C).
\end{equation}

GCM 同时提供机密性和完整性，其底层 AES 标准由 NIST FIPS 197 定义\cite{fips197}。根密钥由部署环境的 Secret Manager/KMS 注入，Agent 数据库只保存密文、keyId 和版本。密钥轮换采用双读单写：新写入使用当前 keyId，读取可使用当前与上一版本，后台重加密完成后撤销旧版本。API 只返回 \code{hasApiKey} 或 \code{hasSecret}，日志脱敏器会移除 Token、Key、密码、私钥和带凭据 URL。

生产入口由反向代理终止 TLS，浏览器只访问同源 HTTPS。CORS 使用显式 Origin 列表，内部端口由网络策略隔离。Agent 对外模型端点使用协议与主机允许列表，远程 MCP 插件使用独立认证头；SSH 出口只允许登记主机与端口。ClickHouse、PostgreSQL、Agent SQLite、Query Internal API 和管理端点不直接暴露到公共网络。

\subsection{审计与追溯}

审计覆盖登录主体、模型设置、工具策略、SSH 主机、命令定义、Loop 定义、调度触发、Run 状态、工具调用、审批和密钥轮换。每条审计事件包含 eventId、orgId、actorType、actorId、action、resourceType、resourceId、result、reason、requestId、runId、traceId、sourceIp、userAgent 和 timestamp。高敏感字段只保存哈希或摘要。

Run Event 与安全审计使用只追加写入。管理界面支持按用户、资源、动作、结果、时间和 traceId 查询；导出任务按时间窗口生成 JSON Lines，并包含完整性哈希。审计保留周期与遥测 TTL 独立配置，删除业务对象不会级联删除审计记录。SSH 输出正文不进入审计表，只保存字节数、内容哈希、截断标记和错误摘要；需要查看执行证据时从具有权限的原 Run 读取。

\begin{table*}[t]
\caption{安全控制与强制验证点}
\label{tab:security-controls}
\centering
\small
\begin{tabularx}{\textwidth}{p{2.3cm}Y Y}
\toprule
\textbf{控制域} & \textbf{执行路径} & \textbf{审计结果} \\
\midrule
遥测访问 & 用户 JWT 透传 Query Service；RBAC 与 orgId 双重过滤；Agent 不持有 ClickHouse 凭据 & 记录用户、查询工具、时间窗口、结果规模和 traceId \\
Prompt 数据边界 & 系统/管理员/用户/附件/工具分层；附件类型与大小校验；敏感字段遮蔽 & 记录策略版本、输入摘要和阻断原因 \\
工具授权 & Zod Schema、角色、资源所有权、风险等级、审批模式和结果上限联合判定 & 记录 toolCallId、参数摘要、决策和耗时 \\
SSH 执行 & 指纹固定、固定命令 ID、无 PTY、无交互、超时、输出上限、逐次审批 & 记录主机、命令、策略版本、审批人和退出状态 \\
秘密保护 & AES-256-GCM、KMS Key ID、双读单写轮换、响应与日志脱敏 & 记录创建、使用、轮换和撤销事件，不记录明文 \\
事件入口 & 服务身份、请求签名、时间窗口、eventId 幂等和载荷 Schema & 记录 source、eventType、匹配结果和 Run 关联 \\
资源治理 & 用户/组织/Loop 并发、速率、Token、费用和输出预算 & 记录配额维度、当前值、阈值和拒绝码 \\
\bottomrule
\end{tabularx}
\end{table*}

\section{可靠性、性能与可运维性}

\subsection{启动顺序与健康模型}

服务依赖形成有向启动图：ZooKeeper $\rightarrow$ ClickHouse $\rightarrow$ Schema Migrator $\rightarrow$ Query Service $\rightarrow$ Agent/Alertmanager/Collector $\rightarrow$ Frontend。Schema Migrator 在 ClickHouse 就绪后执行结构升级，成功退出后 Query Service 才开始提供完整查询。Agent Service 启动时依次完成配置校验、数据库迁移、密钥可用性检查、Repository 初始化、调度租约注册和 HTTP 监听。

每个服务提供 liveness 和 readiness。liveness 只判断进程事件循环或主线程；readiness 同时检查关键依赖、迁移版本和任务接收能力。Agent readiness 检查 PostgreSQL、WASM SQLite、Query Service 身份端点和内置 MCP Server；外部模型供应商不作为启动前置条件，连接状态由模型测试与运行错误独立呈现。优雅下线时服务停止接收新请求，等待活动流和工具结束，再释放 MCP 与数据库连接。

\subsection{容量边界与资源预算}

表~\ref{tab:limits} 给出默认边界。单请求上限与租户预算同时生效，取更严格值。预算器在模型调用前预留 Token 和费用，在流结束后按实际 usage 结算；预留失败会在模型调用之前返回 429，避免执行到一半才终止。

\begin{table*}[t]
\caption{Agent、Loop 与租户资源边界}
\label{tab:limits}
\centering
\small
\begin{tabularx}{\textwidth}{p{3.0cm}p{3.2cm}Y}
\toprule
\textbf{对象} & \textbf{默认值/范围} & \textbf{执行方式} \\
\midrule
模型输出 Token & 默认 8,192；范围 64--128,000 & 供应商参数与组织预算共同约束 \\
模型步骤与重试 & 8 步；1 次供应商重试 & 达到上限后生成明确终态，不重复副作用工具 \\
模型总体超时 & 120 秒；范围 10--1,800 秒 & Abort Signal 传播到模型流和只读工具 \\
记忆窗口 & 默认 30；范围 1--200 条 & 按 userId/threadId 读取最近消息 \\
TopoMem 候选 & 最多扫描 800 条；默认返回 5 条 & 按组织和作用域过滤后执行路径与分块评分 \\
TopoMem 路径 & 错误率阈值 0.10；最多 3 条；最多 6 个服务节点 & 当前证据不足时显式进入 lexical-fallback \\
记忆正文与分块 & 单条 120,000 字节；单块 1,800 字符；最多 80 块 & 稳定序列化、秘密遮蔽、哈希去重和类型标注 \\
工具结果 & 200,000 字节 & 保留结构、摘要、样本和截断标记 \\
消息附件 & 5 个；单文件 10 MB；文本 512 KB & MIME、扩展名、大小和文件名校验 \\
Loop 结构 & 100 节点；5 层；6 并行分支 & 保存和编译阶段双重校验 \\
循环与 foreach & 50 次/项；并发 1--10 & 每次迭代写入 itemKey 与步骤事件 \\
SSH & 5--120 秒；输出 200 KB & 禁止交互；超时主动关闭连接 \\
用户交互 Run & 默认每用户 3 个活动 Run & 超额请求进入队列或返回配额状态 \\
组织 Run & 默认 20 个活动 Run & Dispatcher 公平调度并按优先级出队 \\
API 速率 & 用户 60 req/min；内部事件 600 req/min & 令牌桶；返回 Retry-After \\
Token/费用预算 & 按日和按月配置 & 预留、结算、告警和硬停止四阶段 \\
\bottomrule
\end{tabularx}
\end{table*}

Query Service 对时间范围、步长、序列数、Limit 和查询超时实施独立保护。前端图表根据可视宽度计算目标点数，Agent 默认使用更窄时间窗口和聚合查询。ClickHouse 资源组区分交互查询、告警评估和后台任务，避免大规模仪表盘刷新阻塞告警与 Agent 取证。

\subsection{高可用与横向扩展}

Frontend、Query Service、Collector 和 Agent Service 均采用无会话或外置状态设计，可在负载均衡器后横向扩展。浏览器流式请求使用连接亲和或支持长连接的网关；Mastra Memory 和运行数据保存在 PostgreSQL，Agent 业务配置与审计保存在可复制备份的 SQLite 文件中。多副本部署共享 PostgreSQL，并通过统一入口保持流连接稳定。

ClickHouse 使用分片与副本扩展写入和查询，ZooKeeper 保存协调元数据；Collector 按 OTLP 连接和队列水平扩展。Query Service 查询缓存使用共享缓存层或稳定的本地缓存失效策略。Alertmanager 以集群模式维护告警去重和通知状态。所有有状态组件部署反亲和规则与持久卷，避免单节点故障同时损坏主副本。

Run Worker 领取任务时写入 workerId、leaseEpoch 和 leaseExpiresAt，并周期续租。Worker 失联后，Recovery Worker 只回收租约过期且没有完成事件的 Run；已完成步骤依据检查点跳过。这个机制使模型调用或查询节点在实例退出后可恢复，同时避免活动 Worker 被第二个副本并发推进。

\subsection{自观测}

Scry 使用 OpenTelemetry 观测自身。HTTP 请求、模型调用、Query Service 工具、Loop 节点、调度触发、事件消费、审批等待和 SSH 执行都创建 Span，并传播 \code{traceparent}。核心属性包括 orgId 的哈希、userId 的哈希、threadId、runId、workflowId、stepId、toolId、provider、model、triggerType、status 和 error.type；Prompt、秘密和完整工具结果不写入 Span。

Agent 暴露请求量、错误率、流持续时间、首 Token 延迟、模型 Token、估算费用、工具耗时、队列深度、Run 状态、步骤重试、审批等待、Cron 延迟、事件 Inbox Lag、幂等冲突、租约切换、SSH 结果和数据库事务耗时等指标。关键告警包括：readiness 连续失败、队列持续增长、调度延迟超过阈值、事件 Dead Letter 增长、模型错误率升高、配额拒绝突增、审批超时积压、SSH 指纹不匹配和密钥解密失败。

结构化日志使用 requestId、traceId 和 runId 关联。日志级别按操作结果划分：正常状态转换使用 info，可恢复依赖错误使用 warn，终态系统错误使用 error，调度扫描和状态细节使用 debug。日志输出经过统一脱敏，并限制对象深度和字符串长度。

\subsection{一致性与恢复}

遥测写入采用最终一致性，Collector 批处理和 ClickHouse 物化视图会引入秒级可见延迟；Query Service 在响应元数据中返回实际时间范围和查询完成时间。平台配置、Loop 定义和运行状态采用事务一致性。Workflow 编辑使用版本号进行乐观锁，客户端提交旧版本时返回 409；Run 始终引用启动时的版本。

对话流使用客户端分支与持久化分支并行消费。客户端断开不会中止持久化分支，除非用户显式点击停止；停止操作写入 cancelled 事件并向模型传播 Abort Signal。Assistant Message 使用 Upsert，网络重连和重复完成回调不会产生两条消息。审批决策和 Trigger Claim 使用唯一键，重复请求返回现有结果。

\section{开发、构建与生产部署}

\subsection{开发环境}

开发环境通过 \code{Makefile.dev} 封装 Compose，保留与生产相同的服务边界。常用命令如下：

\begin{lstlisting}[language=bash]
make -f Makefile.dev dev
make -f Makefile.dev dev-ps
make -f Makefile.dev dev-logs
make -f Makefile.dev dev-backend-restart
make -f Makefile.dev dev-agent-restart
make -f Makefile.dev dev-frontend-restart
make -f Makefile.dev dev-down
\end{lstlisting}

前端使用 Webpack Watch，Agent 使用 \code{tsx watch}，Go 服务通过重启触发编译。环境文件只保存本地非生产配置，模型 Key 和 SSH Secret 通过管理界面写入加密存储。开发入口绑定回环地址，数据卷在 \code{dev-down} 后保留，清理卷使用单独的显式命令。Agent 启动前由一次性 \code{agent-data-init} 任务修正持久卷属主；该任务以只读根文件系统运行，删除全部 Linux Capability 后仅恢复 \code{CHOWN}，完成即退出。Agent 主容器始终使用非 root 用户、\code{cap\_drop=ALL} 与 \code{no-new-privileges}，卷升级不需要放宽主进程权限。

\subsection{端口与网络分区}

默认服务端口为：Web Console 3301、Query Public API 8080、Query Internal API 8085、Agent API 4111、ClickHouse Native/HTTP 9000/8123、Alertmanager 9093、OTLP gRPC/HTTP 4317/4318。生产环境只公开 HTTPS 443 和按需公开 OTLP 接入端口；其他端口位于应用、数据和管理三个网络分区。

应用分区包含 Frontend、Query Service 和 Agent Service；数据分区包含 ClickHouse、ZooKeeper、PostgreSQL 和 Agent SQLite 备份；管理分区包含内部 API、迁移任务和备份任务。网络策略允许 Collector 写 ClickHouse、Query Service 读 ClickHouse、Agent 调 Query Service、MCP 与模型端点、Alertmanager 回调内部 API。浏览器不能直接访问数据分区，Agent 不能直接访问 ClickHouse。

\subsection{生产部署流程}

生产发布使用版本化镜像和不可变配置。部署顺序为：创建命名空间与网络策略，注入数据库和密钥 Secret，部署 ZooKeeper/ClickHouse，执行 Schema Migrator，部署 Query Service 与 Alertmanager，部署 Agent Service，部署 Collector，最后发布 Frontend 与入口网关。每一阶段都等待 readiness，再进入下一阶段。

Agent Service 启动前必须提供主密钥 Key ID、数据库 URL、Query Service URL、允许的模型端点和内部服务令牌。模型业务配置由管理员在系统内管理，不写入镜像。SSH 主机通过管理页面逐台纳管并完成指纹核验。Loop 在草稿状态完成保存和试运行，启用后调度器才注册 Cron 与事件触发器。

滚动发布时，新副本完成迁移兼容性检查并进入 readiness，旧副本撤销 readiness 后停止领取 Run。调度租约自动转移，活动 Run 由旧副本完成或保存检查点后交由新副本恢复。入口网关对 UI Message Stream 设置足够长的空闲超时，并关闭响应缓冲。

\subsection{麒麟 V11 与 LoongArch64 原生部署}

\filepath{deploy/kylin-loong64} 提供完整的宿主机原生部署链。\code{build-node.sh} 下载 Node.js 22 源码与官方 SHA-256 清单，在目标机编译后验证 \code{process.arch=loong64}；\code{build-query.sh} 使用目标机 Go \code{GOARCH=loong64} 和 \code{CGO\_ENABLED=1} 构建 Query Service，并通过 \code{file} 检查二进制架构。Agent 先用 TypeScript 生成 \filepath{dist}，运行阶段只安装生产依赖，不使用 \code{tsx}/\code{esbuild}。

Agent 业务数据库由 \code{sql.js} WebAssembly 驱动，使用与 SQLite 兼容的文件格式、串行事务和临时文件原子重命名落盘；Mastra Memory 和运行存储使用麒麟原生 PostgreSQL。该组合不加载仅提供 x64/arm64 预编译包的 Node 原生 SQLite 扩展。Frontend 构建为静态资源，由 Nginx 同源代理 Query Service 和 Agent Service。systemd 单元统一使用无登录的 \code{scry} 用户，启用 \code{NoNewPrivileges}、\code{ProtectSystem=strict}、\code{ProtectHome} 和受限写目录。

部署流程由 \code{preflight.sh}、\code{setup-postgres.sh}、构建脚本、\code{install-native.sh} 与 \code{acceptance.sh} 组成。预检验证麒麟 V11、LoongArch64、Go/Node/GCC/PostgreSQL/Nginx，并通过原生客户端对 ClickHouse 执行 \code{SELECT 1}。验收脚本检查 systemd 状态、HTTP 健康接口、服务用户、数据目录属主和可选 MCP 工具目录，从而把平台适配从说明性声明变为可重复执行的交付步骤。

\subsection{备份与恢复}

备份分为遥测数据、平台配置和智能执行状态三个集合。ClickHouse 使用分区备份或对象存储备份；Query Service SQLite/数据库使用一致性备份；Agent 数据库在事务边界生成备份并单独保存。KMS 根密钥不进入数据库备份，Key ID、密钥版本和恢复权限通过 Secret Manager 的灾备机制保存。

默认策略为每日全量、每 15 分钟增量或日志归档，保留 7 个日版本、4 个周版本和 12 个月版本。生产目标为配置与 Agent 状态 RPO 15 分钟、RTO 60 分钟；遥测数据的 RPO/RTO 根据保留量和 ClickHouse 规模单独计算。备份任务完成后执行可读性校验、对象计数、校验和与密钥版本检查。

恢复顺序为：恢复密钥访问，恢复 ZooKeeper/ClickHouse 与模式，恢复 Query Service 配置库，恢复 Agent 库，执行迁移检查，启动 Query Service，启动 Agent，最后启动 Collector 和 Frontend。Agent 恢复后，Recovery Worker 扫描活动 Run；超过审批有效期的 suspended Run 标记 expired，租约过期的 running Run 重新入队，已完成 Run 保持终态。每次恢复演练记录恢复点、耗时和对象数量。

\subsection{升级与数据库迁移}

所有服务版本遵循向前兼容的滚动升级顺序。数据库迁移采用 expand-and-contract：先添加新表、新列和索引，使新旧版本都可读写；完成应用滚动升级和数据回填后，再在后续版本移除旧结构。迁移表记录 version、checksum、startedAt、completedAt 和 executorId，多副本通过迁移锁保证只执行一次。

Workflow DSL 和 A2UI 都携带版本号。编译器保留对已发布 DSL 版本的读取能力，保存时可将旧定义转换为当前版本并生成新的 Workflow Version；运行中的 Run 继续使用原版本。A2UI 前端对已知版本进行专用渲染，对更高版本保留 JSON 回退。模型供应商适配层以协议能力检测控制参数，升级模型不会改变已开始 Run 的配置副本。

\subsection{运维 Runbook}

日常运维以健康、队列、依赖、配额和运行状态为主线。

\begin{enumerate}
  \item \textbf{Agent 无法访问：}检查入口 5xx、Agent readiness、PostgreSQL、WASM SQLite 和 Query Service 身份端点；使用 requestId/traceId 定位认证、MCP 或依赖错误。
  \item \textbf{模型响应异常：}在 AI 设置执行连接测试，核对 provider、model、endpoint、超时、限流和组织预算；供应商错误通过 category 分类。
  \item \textbf{Loop 长时间排队：}检查实例/组织/Loop 并发、Dispatcher 队列深度、Worker 租约和 nextAttemptAt；扩容 Worker 后队列自动消化。
  \item \textbf{Cron 没有按时运行：}检查 Loop enabled、Cron、timezone、Scheduler Leader、nextFireAt、misfire 记录和 Trigger Claim；按 triggerKey 判断是否被幂等抑制。
  \item \textbf{事件没有触发：}按 eventId 查询 Inbox，检查签名、eventType、filter、orgId、matched 状态和 Dead Letter；修复后使用 replay 端点重投。
  \item \textbf{审批无法恢复：}检查 Run 是否 suspended、审批是否到期、当前用户角色、suspendedStep 和 workflowVersion；409 表示状态已被其他决策推进。
  \item \textbf{SSH 连接失败：}依次检查主机启用状态、网络、指纹、用户名、凭据、策略和超时；指纹变化必须重新走可信核验。
  \item \textbf{查询结果过大：}缩小时间范围、增加步长、减少 Group By 和 Limit；工具返回中的 truncated 与 originalBytes 用于确认截断。
  \item \textbf{恢复后 Run 异常：}查询 Run Event 的最后序号和 worker lease，确认 Recovery Worker 是否重新入队；副作用步骤依据 executionKey 核查审计后再处理。
\end{enumerate}

\section{兼容性与扩展机制}

\subsection{模型与工具扩展}

模型适配器以统一的 generate/stream/tool-call 接口隔离 OpenAI Responses 与 Anthropic Messages。新增供应商时实现模型工厂、Endpoint 规范化、错误分类、连接测试和 usage 映射，不改变 Agent、Thread、A2UI 或 Loop。模型设置保存 provider、protocol、model 和参数副本，运行事件记录实际使用值。

新增工具需要声明唯一 toolId、输入/输出 Schema、风险等级、角色、审批模式、超时、结果上限、A2UI 映射和审计字段。只读工具通过当前用户 JWT 访问领域服务；有副作用工具实现 executionKey 和明确终态。工具注册完成后即可被普通 Agent 和 Loop Agent 节点复用，前端对标准 A2UI 类型无需新增页面逻辑。

\subsection{Loop 与事件扩展}

Loop DSL 通过 version 和节点 type 扩展。新节点实现验证器、编译器、运行事件映射、编辑器配置和迁移函数；旧版本定义保持可执行。事件系统使用开放的 eventType 命名空间，新增事件源只需生成标准信封并配置服务身份。过滤器只读取信封，不与具体来源代码耦合。

Webhook、工单、变更平台和配置中心等外部系统通过 Integration Adapter 转换为内部事件。出站动作通过受控工具实现，继续复用角色、审批、幂等、预算和审计，不在 Loop 编译器中直接写第三方 API 逻辑。

\section{结语}

Scry 将遥测接入、分析查询、平台管理和智能执行组织为边界清晰的四平面架构。指标、日志与链路由 OpenTelemetry 和 ClickHouse 提供统一事实基础；Query Service 集中承担查询语义、RBAC、告警、仪表盘和配置；Agent Service 则以 TopoMem、证据链、确定性 MCP 和安全 Loop 形成系统的技术主线。四者分别解决“先查哪里”“依据什么”“如何可靠调用”和“怎样重复执行”，使智能诊断不脱离可观测平台的真实数据、权限和运行状态。

TopoMem 使用本次事故 Trace 导航到异常调用链下游，以错误签名和记忆分块区分同一服务的不同故障模式，并在 Trace 缺失、终端无记忆或多路径无支持时显式回退。它与证据系统共同构成验证后学习闭环：历史记忆只生成假设，本轮工具结果形成证据，人工接受或 Loop 验证后的结果才能写回，维护器再执行衰减、强化、合并、冲突挂起和归档。MCP 与 Loop 位于闭环的执行侧，前者提供严格参数与权限边界，后者提供版本、审批、检查点、幂等和恢复语义。由此，Scry 把一次自然语言诊断转化为可解释、可治理、可积累并可重复运行的技术过程。

\clearpage
\nobalance
\onecolumn
\appendices

\section{Agent API 参考}

\begin{table}[htbp]
\caption{Agent Service HTTP API}
\label{tab:appendix-api}
\centering
\scriptsize
\begin{tabularx}{\textwidth}{p{1.6cm}p{5.8cm}p{2.1cm}Y}
\toprule
\textbf{方法} & \textbf{路径} & \textbf{权限} & \textbf{语义} \\
\midrule
GET & \code{/health} & Public & 进程、版本、readiness 摘要 \\
GET & \code{/api/settings/model} & ADMIN & 返回模型配置和秘密存在状态 \\
PUT & \code{/api/settings/model} & ADMIN & 更新供应商、端点、模型、参数和记忆窗口 \\
POST & \code{/api/settings/model/test} & ADMIN & 使用表单与已保存秘密合并值执行连接测试 \\
POST & \code{/mcp} & Login & MCP Streamable HTTP initialize、tools/list 与 tools/call \\
GET & \code{/api/mcp/tools} & Login & 返回当前用户可见的内置 MCP 工具目录 \\
GET & \code{/api/mcp/servers} & ADMIN & 按组织列出远程 MCP 插件，不返回认证头明文 \\
POST & \code{/api/mcp/servers} & ADMIN & 注册远程 Streamable HTTP MCP 服务 \\
PUT & \code{/api/mcp/servers/:serverId} & ADMIN & 更新插件、超时、启用状态和加密请求头 \\
DELETE & \code{/api/mcp/servers/:serverId} & ADMIN & 删除组织内 MCP 插件 \\
POST & \code{/api/mcp/servers/:serverId/test} & ADMIN & 建立 MCP 连接并列出工具 \\
GET & \code{/api/ssh/hosts} & ADMIN & 列出组织内 SSH 主机，不返回秘密 \\
POST & \code{/api/ssh/hosts/probe} & ADMIN & 读取远端主机公钥指纹 \\
POST & \code{/api/ssh/hosts} & ADMIN & 创建主机并加密保存认证材料 \\
GET & \code{/api/ssh/hosts/:hostId} & ADMIN & 获取主机配置和秘密存在状态 \\
PUT & \code{/api/ssh/hosts/:hostId} & ADMIN & 更新主机、指纹、凭据和启用状态 \\
DELETE & \code{/api/ssh/hosts/:hostId} & ADMIN & 删除主机与策略引用 \\
POST & \code{/api/ssh/hosts/:hostId/test} & ADMIN & 执行指纹与认证握手测试 \\
GET & \code{/api/ssh/policies} & ADMIN & 返回只读与扩展执行策略 \\
PUT & \code{/api/ssh/policies/:toolId} & ADMIN & 更新启用、审批、超时和命令定义 \\
GET & \code{/api/ssh/audits} & ADMIN & 按用户、主机、工具、结果和时间查询审计 \\
GET & \code{/api/security/audits} & ADMIN & 按 traceId 查询输入、策略、审批和执行事件 \\
GET & \code{/api/threads} & Login & 按当前用户列出 Thread \\
POST & \code{/api/threads} & Login & 创建 Thread \\
GET & \code{/api/threads/:threadId} & Owner & 返回 Thread 与完整 UI Message \\
DELETE & \code{/api/threads/:threadId} & Owner & 删除 Thread、消息和关联运行 \\
POST & \code{/api/chat} & Owner & 普通 Agent 流式请求、重新生成或恢复 \\
GET & \code{/api/threads/:threadId/evidence} & Owner & 返回证据并从历史工具 Part 补建缺失证据 \\
PATCH & \code{/api/evidence/:evidenceId} & Owner & 更新接受、驳回、挂起、选择状态或备注 \\
POST & \code{/api/threads/:threadId/evidence/revise} & Owner & 使用选定证据生成回答修订并记录排除项 \\
GET & \code{/api/memories/stats} & Login & 返回总量、在线量、验证量、挂起量、平均活性和分布 \\
GET & \code{/api/memories} & Login & 按查询、类型、作用域、状态和服务过滤记忆 \\
POST & \code{/api/memories} & Login/Admin & 创建记忆；组织级写入要求 ADMIN；可绑定证据 UUID \\
GET & \code{/api/memories/:memoryId} & Visible & 返回正文、分块、来源、关系和维护历史 \\
PATCH & \code{/api/memories/:memoryId} & Visible/Admin & 更新状态、固定、备注、置信度和重要性 \\
DELETE & \code{/api/memories/:memoryId} & Visible/Admin & 清空可检索内容并保留遗忘审计骨架 \\
POST & \code{/api/memories/search} & Login & 执行 TopoMem/TopoMem-MP/lexical-fallback 并返回评分解释 \\
GET & \code{/api/threads/:threadId/memories} & Owner & 使用当前 Thread 问题与证据生成记忆包 \\
POST & \code{/api/memories/maintenance} & Login & 执行衰减、强化、合并、冲突、归档和过期处理 \\
GET & \code{/api/memory-samples} & Login & 返回按类别组织的论文语料示例及生成状态 \\
POST & \code{/api/memory-samples/generate} & Login & 选择示例并生成为正常参与检索的记忆 \\
DELETE & \code{/api/memory-samples/:memoryId} & Login & 删除指定示例记忆及分块、来源、关系和操作记录 \\
GET & \code{/api/workflows} & Login & 列出有权查看的 Loop \\
POST & \code{/api/workflows} & ADMIN & 创建并验证 Loop 定义 \\
GET & \code{/api/workflows/:workflowId} & Login & 返回定义、输入、触发器、版本和下次执行 \\
PUT & \code{/api/workflows/:workflowId} & ADMIN & 乐观锁更新并生成版本记录 \\
DELETE & \code{/api/workflows/:workflowId} & ADMIN & 停用调度并删除定义 \\
GET & \code{/api/workflows/:workflowId/runs} & Login & 查询当前主体可见的运行历史 \\
POST & \code{/api/workflows/:workflowId/stream} & Login & 手动创建 Run 并返回 UI Message Stream \\
POST & \code{/api/workflows/:workflowId/trigger} & ADMIN & 通过统一触发管线创建手动 Run \\
GET & \code{/api/workflow-runs/:runId} & Owner/Admin & 返回 Run、步骤、触发器、统计和输出 \\
GET & \code{/api/workflow-runs/:runId/events} & Owner/Admin & 返回有序 Run Event \\
POST & \code{/api/workflow-runs/:runId/resume} & Approver & 对 suspended Run 提交批准或拒绝 \\
POST & \code{/api/workflow-runs/:runId/cancel} & Owner/Admin & 取消 queued/running/retrying/suspended Run \\
GET & \code{/api/scheduler/status} & ADMIN & 返回 Leader、租约、队列、延迟和任务注册 \\
POST & \code{/api/internal/events} & Service & 验证并写入事件信封 \\
POST & \code{/api/internal/events/:eventId/replay} & ADMIN & 重新投递 Dead Letter 或历史事件 \\
\bottomrule
\end{tabularx}
\end{table}

\clearpage
\section{Agent 数据模型}

\begin{table}[htbp]
\caption{Agent 关键表与字段}
\label{tab:appendix-data-model}
\centering
\scriptsize
\begin{tabularx}{\textwidth}{p{4.0cm}p{5.2cm}Y}
\toprule
\textbf{表} & \textbf{关键字段} & \textbf{职责与约束} \\
\midrule
\filepath{scry_model_settings} & orgId, environment, provider, endpointMode, model, baseUrl, apiKeyEncrypted, keyVersion, generationOptions, updatedAt & 组织级模型配置；唯一键为 orgId+environment；API 只返回秘密存在状态 \\
\filepath{scry_agent_threads} & id, orgId, userId, title, createdAt, updatedAt & Thread 所有权与列表排序 \\
\filepath{scry_agent_ui_messages} & id, threadId, role, messageJson, createdAt, updatedAt & 保存完整 UI Message Part；messageId Upsert \\
\filepath{scry_evidence_items} & id, threadId, orgId, userId, traceId, messageId, sourceType, sourceName, contentJson, confidence, relevance, status, selected, note, capturedAt & 工具与文档证据主体；保存评分、生命周期和用户选择 \\
\filepath{scry_evidence_relations/actions/revisions} & evidenceId, relationType, score, action, previousStatus, nextStatus, selectedIdsJson, excludedIdsJson, instruction, createdAt & 证据关系、人工处置和基于更新证据的回答修订 \\
\filepath{scry_memory_items} & id, orgId, userId, threadId, workflowId, memoryType, scope, serviceName, contentHash, confidence, importance, vitality, status, pinned, expiresAt & TopoMem 主记录；按作用域、服务、状态和活性检索 \\
\filepath{scry_memory_chunks} & id, memoryId, chunkIndex, chunkType, contentText, contentHash, tokenCount, keywordsJson, serviceName, confidence & 稳定分块、预览和块级检索；memoryId+chunkIndex 唯一 \\
\filepath{scry_memory_sources/relations/actions} & memoryId, sourceType, sourceId, traceId, relation, weight, relationType, score, action, status, detail & 记忆来源、支持/冲突/重复关系及维护历史 \\
\filepath{scry_memory_topology_edges} & orgId, fromService, toService, observationCount, errorObservationCount, cumulativeErrorRate, confidence, lastTraceId, lastObservedAt & 运行拓扑观察累计；不替代当前事故 Trace \\
\filepath{scry_memory_retrieval_runs} & id, orgId, userId, threadId, traceId, queryText, strategy, symptomService, signatureJson, pathsJson, fallbackReason, candidateCount, returnedIdsJson, latencyMs & TopoMem 检索决策审计和可解释复盘 \\
\filepath{scry_mcp_servers} & id, orgId, name, url, enabled, headersEncrypted, timeoutMs, createdAt, updatedAt & 远程 MCP 插件注册；认证头只以 AES-GCM 密文保存 \\
\filepath{scry_ssh_hosts} & id, orgId, name, hostname, port, username, authType, secretEncrypted, keyVersion, fingerprint, enabled & SSH 主机、认证密文和固定指纹 \\
\filepath{scry_tool_policies} & orgId, toolId, riskLevel, enabled, requireApproval, policyJson, version, updatedAt & 工具启用、命令定义、超时和审批版本 \\
\filepath{scry_ssh_audit_logs} & id, orgId, userId, runId, toolCallId, hostId, commandId, policyVersion, status, exitCode, durationMs, stdoutBytes, stderrBytes, outputHash, createdAt & SSH 成功、失败和拒绝审计 \\
\filepath{scry_security_audit_events} & id, traceId, orgId, userId, eventType, source, target, decision, riskScore, inputHash, detailsJson, createdAt & 输入扫描、意图分类、策略、审批和工具执行统一审计 \\
\filepath{scry_workflows} & id, orgId, ownerUserId, name, description, enabled, definitionJson, tagsJson, scheduleJson, eventTriggersJson, inputSchemaJson, version, concurrencyPolicy, executionPrincipal, createdAt, updatedAt & Loop 当前版本、触发器、身份和并发策略 \\
\filepath{scry_workflow_versions} & workflowId, version, definitionJson, checksum, createdBy, createdAt & 不可变历史版本；workflowId+version 唯一 \\
\filepath{scry_workflow_runs} & id, orgId, userId, workflowId, workflowVersion, threadId, triggerType, triggerKey, status, attempt, inputJson, outputJson, stepStateJson, suspendedStep, scheduledAt, nextAttemptAt, workerId, leaseEpoch, createdAt, updatedAt & Run 当前状态副本；状态更新使用 expectedStatus \\
\filepath{scry_workflow_run_events} & runId, sequence, eventType, actor, payloadSummary, traceId, createdAt & 只追加状态事件；runId+sequence 唯一 \\
\filepath{scry_trigger_claims} & triggerKey, orgId, workflowId, triggerType, sourceId, scheduledAt, runId, createdAt & Cron/Event 幂等；triggerKey 唯一 \\
\filepath{scry_event_inbox} & eventId, replayId, orgId, eventType, source, envelopeJson, status, attempts, nextAttemptAt, errorCode, receivedAt, processedAt & 事件持久接收、匹配、重试与 Dead Letter \\
\filepath{scry_runtime_leases} & name, holderId, epoch, expiresAt, updatedAt & Scheduler Leader 与 Worker Run Lease \\
\filepath{scry_agent_audit_events} & eventId, orgId, actorType, actorId, action, resourceType, resourceId, result, requestId, runId, traceId, metadataJson, createdAt & Agent 管理和执行统一审计 \\
\bottomrule
\end{tabularx}
\end{table}

\clearpage
\section{Agent 配置环境变量}

\begin{table}[htbp]
\caption{Agent Service 环境变量}
\label{tab:appendix-env}
\centering
\scriptsize
\begin{tabularx}{\textwidth}{p{5.8cm}p{3.0cm}Y}
\toprule
\textbf{变量} & \textbf{默认值} & \textbf{含义} \\
\midrule
\filepath{SCRY_AGENT_PORT} & 4111 & HTTP 监听端口 \\
\filepath{SCRY_AGENT_HOST} & \code{0.0.0.0} & 容器内监听地址 \\
\filepath{SCRY_QUERY_SERVICE_URL} & 必填 & Query Service Public API 地址 \\
\filepath{SCRY_QUERY_INTERNAL_URL} & 必填 & 身份换取与内部事件协作地址 \\
\filepath{SCRY_AGENT_DATABASE_URL} & \code{file:/var/lib/scry/agent/agent.db} & WASM SQLite 业务数据库文件 \\
\filepath{SCRY_AGENT_POSTGRES_URL} & 必填 & Mastra Memory 与运行存储的 PostgreSQL URL \\
\filepath{SCRY_MCP_INTERNAL_URL} & \code{http://127.0.0.1:4111/mcp} & 内置 MCP Server 地址 \\
\filepath{SCRY_MCP_TIMEOUT_MS} & 30000 & MCP 初始化、工具目录和调用超时 \\
\filepath{SCRY_AGENT_MASTER_KEY} & 必填 & 本地 Secret 模式下的高熵主密钥 \\
\filepath{SCRY_AGENT_KMS_KEY_ID} & 生产必填 & KMS 当前加密 Key ID \\
\filepath{SCRY_AGENT_PREVIOUS_KMS_KEY_ID} & 空 & 轮换期间允许读取的上一 Key ID \\
\filepath{SCRY_AGENT_MAX_TOOL_RESULT_BYTES} & 200000 & 单次工具结果上限 \\
\filepath{SCRY_AGENT_MAX_ACTIVE_RUNS} & 100 & 实例活动 Run 上限 \\
\filepath{SCRY_AGENT_MAX_ORG_RUNS} & 20 & 单组织活动 Run 上限 \\
\filepath{SCRY_AGENT_MAX_USER_RUNS} & 3 & 单用户交互 Run 上限 \\
\filepath{SCRY_AGENT_RUN_LEASE_SECONDS} & 30 & Worker Run Lease 时长 \\
\filepath{SCRY_AGENT_SCHEDULER_LEASE_SECONDS} & 15 & Scheduler Leader Lease 时长 \\
\filepath{SCRY_AGENT_SCHEDULER_SCAN_SECONDS} & 5 & Cron 到期扫描周期 \\
\filepath{SCRY_AGENT_EVENT_BATCH_SIZE} & 100 & Event Inbox 单批消费数量 \\
\filepath{SCRY_AGENT_EVENT_MAX_ATTEMPTS} & 8 & 事件消费最大尝试次数 \\
\filepath{SCRY_AGENT_EVENT_SIGNATURE_SECRET} & 必填 & 内部事件 HMAC 服务密钥 \\
\filepath{SCRY_AGENT_ALLOWED_MODEL_HOSTS} & 官方域名 & 自定义模型端点允许列表 \\
\filepath{SCRY_AGENT_CORS_ORIGINS} & 同源 & 允许访问 Agent API 的 Origin 列表 \\
\filepath{SCRY_AGENT_OTEL_ENDPOINT} & Collector 地址 & OTLP Span、Metric 和 Log 导出端点 \\
\filepath{SCRY_AGENT_LOG_LEVEL} & \code{info} & 结构化日志级别 \\
\filepath{SCRY_AGENT_SHUTDOWN_GRACE_SECONDS} & 30 & 优雅下线等待时间 \\
\bottomrule
\end{tabularx}
\end{table}

\clearpage
\section{Loop 定义示例}

下面的定义表示：读取服务与告警，并行生成诊断和仪表盘上下文；当存在严重告警时进入人工审批，批准后执行 Agent 处置步骤。该 Loop 每 15 分钟运行，也响应 \code{alert.firing} 事件。

\begin{lstlisting}
{
  "version": 2,
  "name": "production-health-diagnosis",
  "enabled": true,
  "tags": ["production", "health"],
  "schedule": {
    "enabled": true,
    "cron": "*/15 * * * *",
    "timezone": "Asia/Shanghai",
    "misfirePolicy": "fire-once"
  },
  "eventTriggers": [
    {
      "id": "critical-alert",
      "eventType": "alert.firing",
      "filter": "attributes.severity in ['critical','high']",
      "enabled": true
    }
  ],
  "inputSchema": [
    {
      "id": "environment",
      "label": "environment",
      "type": "text",
      "required": true,
      "defaultValue": "production"
    }
  ],
  "concurrencyPolicy": "forbid",
  "definition": {
    "version": 2,
    "steps": [
      { "id": "services", "name": "list-services", "type": "services" },
      { "id": "alerts", "name": "list-alerts", "type": "alerts" },
      {
        "id": "parallel-context",
        "name": "collect-context-in-parallel",
        "type": "parallel",
        "parallel": {
          "branches": [
            {
              "id": "diagnosis",
              "name": "diagnosis",
              "steps": [{
                "id": "analyze",
                "name": "analyze-root-cause",
                "type": "agent",
                "prompt": "Locate the root cause from services and alerts."
              }]
            },
            {
              "id": "dashboards",
              "name": "dashboards",
              "steps": [{
                "id": "dashboard-list",
                "name": "list-dashboards",
                "type": "dashboards"
              }]
            }
          ]
        }
      },
      {
        "id": "critical-check",
        "name": "critical-alert-check",
        "type": "condition",
        "condition": {
          "rule": {
            "path": "results.alerts.criticalCount",
            "operator": "gt",
            "value": "0"
          },
          "whenTrue": [
            {
              "id": "approve",
              "name": "approve-operation",
              "type": "approval",
              "approvalMessage": "Approve the production operation"
            },
            {
              "id": "operate",
              "name": "execute-operation",
              "type": "agent",
              "prompt": "Execute the approved operation and verify recovery."
            }
          ],
          "whenFalse": []
        }
      }
    ]
  }
}
\end{lstlisting}

\section{运行状态与错误分类}

\begin{table}[htbp]
\caption{Run 状态定义}
\label{tab:appendix-run-status}
\centering
\small
\begin{tabularx}{\textwidth}{p{2.8cm}p{4.2cm}Y}
\toprule
\textbf{状态} & \textbf{允许来源} & \textbf{语义与后续状态} \\
\midrule
\code{queued} & Trigger Claim 成功 & 等待 Worker；进入 running 或 cancelled \\
\code{running} & Worker 领取 queued/retrying & 执行步骤；进入 suspended、retrying、completed、failed 或 cancelled \\
\code{suspended} & approval/tool approval & 保存检查点等待决策；进入 queued、failed、expired 或 cancelled \\
\code{retrying} & 可重试依赖或供应商错误 & 保存 nextAttemptAt；到期进入 queued，超过次数进入 failed \\
\code{completed} & 所有步骤成功 & 终态；保存输出、usage、费用和完成事件 \\
\code{failed} & 终态错误或审批拒绝 & 终态；保存 error code、message、traceId 和步骤状态 \\
\code{cancelled} & 用户、管理员或 replace 策略 & 终态；传播 Abort Signal 并记录取消主体 \\
\code{expired} & 审批或运行有效期到期 & 终态；保留恢复点与到期原因 \\
\bottomrule
\end{tabularx}
\end{table}

\begin{table}[htbp]
\caption{API 与运行错误分类}
\label{tab:appendix-errors}
\centering
\small
\begin{tabularx}{\textwidth}{p{3.8cm}p{2.1cm}p{2.1cm}Y}
\toprule
\textbf{错误码} & \textbf{HTTP} & \textbf{Run 处理} & \textbf{含义} \\
\midrule
\code{VALIDATION\_ERROR} & 400 & failed & Schema、Cron、过滤器、DSL 或参数错误 \\
\code{AUTHENTICATION\_REQUIRED} & 401 & failed & Token 缺失、过期或签名无效 \\
\code{FORBIDDEN} & 403 & failed & 角色、组织、所有权或工具策略拒绝 \\
\code{NOT\_FOUND} & 404 & failed & Thread、Workflow、Run、主机或命令不存在 \\
\code{STATE\_CONFLICT} & 409 & 保持当前状态 & 版本、审批、幂等或状态转换冲突 \\
\code{QUOTA\_EXCEEDED} & 429 & queued/failed & 速率、并发、Token 或费用预算达到阈值 \\
\code{PROVIDER\_RATE\_LIMIT} & 429 & retrying & 模型供应商限流 \\
\code{DEPENDENCY\_UNAVAILABLE} & 503 & retrying & Query Service、数据库、模型或 SSH 依赖暂时不可用 \\
\code{DEPENDENCY\_TIMEOUT} & 504 & retrying & 外部依赖超过超时 \\
\code{APPROVAL\_REJECTED} & 409 & failed & 审批人拒绝执行 \\
\code{RUN\_CANCELLED} & 409 & cancelled & 用户、管理员或并发策略取消 \\
\code{INTERNAL\_ERROR} & 500 & failed/retrying & 平台内部异常，响应携带 traceId \\
\bottomrule
\end{tabularx}
\end{table}

\clearpage
\twocolumn
\bibliographystyle{IEEEtran}
\bibliography{references}

\end{document}
